\documentclass[a4paper, 11pt]{article}
\usepackage[utf8]{inputenc}
\usepackage[
  left = 20mm, right = 20mm, top = 22mm, bottom = 22mm,
]{geometry}

\usepackage{amsmath,amsfonts,amssymb,bm}
\usepackage[english]{babel}
\usepackage{natbib}
\usepackage{hyperref}
\usepackage{orcidlink}
\usepackage{xcolor}
\usepackage{fancyhdr}
\usepackage{graphicx}
\usepackage{booktabs}
\usepackage{float}

\definecolor{harvardcrimson}{HTML}{8C1D40}

\hypersetup{
  colorlinks=true,
  linkcolor=harvardcrimson,
  citecolor=harvardcrimson,
  urlcolor=harvardcrimson,
  hypertexnames=false,
}

\pagestyle{fancy}
\fancyhf{}
\rhead{\textcolor{harvardcrimson}{\textbf{\texttt{RichnessSelection}: counting eligible galaxies}}}
\lhead{\textcolor{harvardcrimson}{\textbf{J.\,H.\,Esteves}}}
\cfoot{\thepage}
\renewcommand{\headrulewidth}{0.4pt}
\renewcommand{\footrulewidth}{0.4pt}
\setlength{\headheight}{14pt}

\renewcommand{\baselinestretch}{1.15}
\setlength{\emergencystretch}{3em}
\sloppy

\renewcommand{\topfraction}{0.95}
\renewcommand{\bottomfraction}{0.9}
\renewcommand{\textfraction}{0.05}
\renewcommand{\floatpagefraction}{0.95}
\setlength{\intextsep}{8pt}
\setlength{\textfloatsep}{10pt}
\setlength{\floatsep}{10pt}

\newcommand{\lob}{\lambda^{\rm ob}}
\newcommand{\ltr}{\lambda^{\rm tr}}
\newcommand{\zob}{z^{\rm ob}}
\newcommand{\dprj}{\Delta^{\rm prj}}
\newcommand{\ximm}{\xi_{\rm mm}}
\newcommand{\xihm}{\xi_{\rm hm}}
\newcommand{\hMpc}{h^{-1}\,\mathrm{Mpc}}
\newcommand{\hMsun}{h^{-1}\,M_{\odot}}
\newcommand{\Pop}{\mathcal{P}}
\newcommand{\avg}[1]{\langle #1 \rangle}
\newcommand{\Rex}{R_{\rm excl}}
\newcommand{\rp}{s}
\newcommand{\nlss}{n_{\rm lss}}
\newcommand{\nrnd}{n_{\rm rnd}}
\newcommand{\rhoprj}{\rho^{\rm prj}}
\newcommand{\bsel}{b_{\rm sel}}
\newcommand{\btar}{b_{\rm target}}
\newcommand{\bsmall}{b_{\rm sel}^{\rm small}}
\newcommand{\blarge}{b_{\rm sel}^{\rm large}}
\newcommand{\Ncal}{\mathcal{N}}
\newcommand{\Nbar}{\bar{N}_{\rm rnd}}
\newcommand{\Klss}{K_{\rm lss}}
\newcommand{\KlssInf}{K_{\rm lss}^{\infty}}
\newcommand{\Nlss}{N_{\rm lss}}
\newcommand{\Nrnd}{N_{\rm RND}}
\newcommand{\thob}{\theta_{\lambda,{\rm ob}}}
\newcommand{\brm}{b_{\rm rm}}
\newcommand{\brmss}{b_{\rm rm}^{\rm ss}}
\newcommand{\brmls}{b_{\rm rm}^{\rm ls}}
\newcommand{\bhalo}{b_{\rm halo}}
\newcommand{\Dl}{\langle \Delta\lambda\rangle}
\newcommand{\Dp}{\langle \dprj\rangle}
\newcommand{\tDl}{\langle \widetilde{\Delta\lambda}\rangle}
\newcommand{\tDp}{\langle \widetilde{\Delta}^{\rm prj}\rangle}
% channel operators (single-sourced; use these, never retype)
\newcommand{\Dlrnd}{\langle \Delta\lambda\rangle_{\rm rnd}}
\newcommand{\Dllss}[1]{\langle \Delta\lambda(#1)\rangle_{\rm lss}}
\newcommand{\tDllss}[1]{\langle \widetilde{\Delta\lambda}(#1)\rangle_{\rm lss}}
\newcommand{\Dprnd}{\langle \dprj\rangle_{\rm rnd}}
\newcommand{\Dpss}{\langle \dprj\rangle^{\rm ss}_{\rm lss}}
\newcommand{\Dpls}{\langle \dprj\rangle^{\rm ls}_{\rm lss}}
\newcommand{\tDpss}{\langle \widetilde{\Delta}^{\rm prj}\rangle^{\rm ss}_{\rm lss}}
\newcommand{\tDpls}{\langle \widetilde{\Delta}^{\rm prj}\rangle^{\rm ls}_{\rm lss}}
\newcommand{\dprjx}{\delta_{\rm prj}}
% section 2 (frozen-algorithm) macros -- bare tokens, external parens,
% e.g. \Ilam(x), \Ilamss(0), \Dlb(z) -- matching the \Dl(z) convention above
\newcommand{\Ilam}{I_\lambda}
\newcommand{\Ilamss}{I_\lambda^{\rm ss}}
\newcommand{\Ilamls}{I_\lambda^{\rm ls}}
\newcommand{\Dlb}{\langle \Delta\lambda\rangle_b}
\newcommand{\pis}{\pi_s}
\newcommand{\rmax}{r_{\max}}
\newcommand{\mulo}{\mu_{\rm lo}}
\newcommand{\muhi}{\mu_{\rm hi}}
\newcommand{\barw}{\bar w}
\newcommand{\Ass}{A_{\rm ss}}
\newcommand{\Bss}{B_{\rm ss}}
\newcommand{\Als}{A_{\rm ls}}
\newcommand{\Bls}{B_{\rm ls}}
\newcommand{\Gss}{G_{\rm ss}}
\newcommand{\DSprj}{\Delta\Sigma^{\rm prj}}
\newcommand{\DSmis}{\Delta\Sigma_{\rm mis}}
\newcommand{\Rmis}{R_{\rm mis}}

\title{\texorpdfstring{%
\textcolor{harvardcrimson}{\textbf{\texttt{RichnessSelection}:\\
counting the eligible galaxies behind the projection integrals}}%
}{RichnessSelection: counting the eligible galaxies behind the projection integrals}}
\author{Johnny H.\ Esteves\,\orcidlink{0000-0003-4373-2386}\\[2pt]
{\small Department of Physics, Harvard University}}
\date{\today}

\begin{document}
\maketitle
\thispagestyle{fancy}

% ------------------------------------------------------------------
% OLD ABSTRACT — commented out during the re-write (kept for reference)
% ------------------------------------------------------------------
\iffalse
\noindent
Projection inflates the observed richness of an optically selected
cluster: $\lob = \ltr + \dprj$, where $\dprj$ counts galaxies
donated by other structures along the line of sight. This note
recasts the projection integrals of \texttt{RichnessSelection}
(\texttt{docs/richness\_selection.tex}, eq.~\texttt{Pop}) as what
they physically are: \emph{weighted counts of eligible galaxies}
in the target's photo-$z$ cylinder, organised by the standard
two-tracer logic --- the target cluster is a biased tracer
($\btar$), every donor halo is a biased tracer ($b(M,z)$), and
their correlation is $\btar\,b(M,z)\,\ximm(r)$. The contamination
then splits exactly like a halo-model profile, one level down:
a Poisson count $\Nbar$ (donors placed at random --- the $b=0$
limit) plus a large-scale-structure count
$\Nlss = \btar\,\Klss$, where $\Klss$ is the donor-bias-weighted
contamination kernel --- one bias short of a complete two-tracer
term, waiting for the target's. Three moves make the counts
nearly analytic. (i)~The contamination halos --- everything
percolation has not already used up --- are contracted once per
$(\lob,\zob)$ into 1-D \emph{counting kernels} $\Ncal_{\rm rnd}(x)$,
$\Ncal_{\rm lss}(x)$: the mean number of galaxies donated into the target
aperture from scaled transverse offset $x$, unweighted and
bias-weighted. (ii)~Separations are decomposed
into a line-of-sight component
$\pi$ and a sky-plane component $\rp$, with 3-D separation
$r^2=\pi^2+\rp^2$ \emph{exactly}; the halo-exclusion ball and the
clustering kernel $\ximm(r)$ become spherical, and the twin peaks
and ring plateau that plague the original $z$-integrand never
form. (iii)~The slowly evolving weights (HMF, bias, MOR, aperture)
are frozen at $\zob$ inside the exclusion-affected zone,
protected by a $10^{3}$ scale separation and by parity. The
Poisson count $\Nbar$ collapses to a 1-D integral by a
total-donation theorem (a donor's disk size redistributes its
galaxies in angle but cannot change their total); the clustered
kernels $\Klss$, $\KlssInf$ become one smooth radial integral plus
one line-of-sight integral. All approximations are validated
against \texttt{scipy.integrate.quad}: the frozen ring alone
costs $0.001$--$0.003\%$, the complete assembly
$0.006$--$0.04\%$ on $(\Nbar, \KlssInf, \Klss)$, and the
end-to-end selection-bias profile $\bsel(\theta)$ agrees with the
production pipeline to $\le0.07\%$ of the profile scale
(\S\ref{sec:bsel_validation}) --- all far below the $0.1\%$
pipeline tolerance.
\fi

\tableofcontents

\section{The physical question: which galaxies contaminate \texorpdfstring{$\lob$}{lambda-ob}?}
\label{sec:question}

\textsc{redMaPPer} estimates a cluster's richness by counting
member galaxies inside a disk of angular radius $\theta_\lambda$
and selecting them with an effective photo-$z$ window around the
cluster, $w(z,\zob)$: the probability that a galaxy at redshift
$z$ passes the photo-$z$ cut of a target at $\zob$ (the parabolic
kernel of the production pipeline; never approximated in this
note). Any galaxy that satisfies both cuts is counted, whether or
not it belongs to the cluster.

\paragraph{Background: a scale-dependent selection bias.}
Richness selection is not mass selection. At fixed mass, systems
whose line of sight carries excess correlated structure scatter
\emph{up} in observed richness, so a sample selected at $\lob$ is
preferentially drawn from overdense environments and clusters
more strongly than a mass-matched sample. The effect is
scale-dependent. At large separations the selected sample behaves
as a linear tracer with a boosted, constant bias amplitude. At
small separations --- within a few richness apertures --- the very
structure that inflated $\lob$ sits inside or near the
measurement itself, and the effective amplitude differs again.
We encode this in a \textsc{redMaPPer} selection bias
$\brm(\theta,\lob,\zob)$, defined as the effective tracer bias of
the selected sample in its cross-correlation with halos of mass
$M$:
\begin{equation}
\xi_{\rm rm,h}(\theta; M, z)
= b(M, z)\; \brm(\theta, \lob, \zob)\;
\ximm\big(|\bm{r}_{\rm rm} - \bm{r}_{\rm h}|\big),
\label{eq:xi_rmh}
\end{equation}
with $b(M,z)$ the halo bias of the correlated structure, $\ximm$
the nonlinear matter--matter correlation function (halofit), and
$|\bm{r}_{\rm rm}-\bm{r}_{\rm h}|$ the 3-D comoving separation
between the target and the halo --- both tracers are biased, and
the cross-correlation carries both biases. Following the
simulation-calibrated model of Costanzi et al., we take
$\brm$ to be two plateaus bridged by a fixed sigmoid,
\begin{equation}
\brm(\theta, \lob, \zob)
= \brmss\,[1-\sigma(\theta)] + \brmls\,\sigma(\theta),
\qquad
\sigma(\theta)=\Big[1+e^{-k(\theta-\theta_0)}\Big]^{-1},
\label{eq:brm}
\end{equation}
with $k=2.5/\thob$ and $\theta_0=\thob/2$, where
$\thob$ is the target's aperture angle: $\sigma$ runs from
$\simeq0$ well inside the aperture to $1$ at large separations,
so $[1-\sigma]$ tags the small-scale regime governed by the
plateau $\brmss$ and $\sigma$ the large-scale regime governed by
$\brmls$. The transition at half the aperture, and the sharpness
$k$, are held fixed; the two plateaus are the physical unknowns.

This bias is not bookkeeping: it propagates directly into the
stacked weak-lensing observable. The selected sample's
large-scale-structure lensing contribution --- the revisited
second-halo term, under the correlation model of
eq.~\eqref{eq:xi_rmh} --- reads
\begin{align}
\avg{\DSprj(R\mid\lob,\zob)} =\;
&2\pi\int d\theta\,\sin\theta
\int dz\,\frac{dV}{dz\,d\Omega}
\int dM\,n(M,z)
\notag\\
&\times\,
\bigl[\,1 + b(M,z)\,\brm(\theta,\lob,\zob)\,
\ximm\big(|\bm{r}_{\rm rm}-\bm{r}_{\rm h}|\big)\,\bigr]
\notag\\
&\times\;
\DSmis\bigl(R\mid M,\Rmis=\theta\,D_A(\zob)\bigr),
\label{eq:DSprj}
\end{align}
where $\DSmis$ is the excess surface density of a halo of mass
$M$ miscentred by the physical transverse offset
$\Rmis = \theta\,D_A(\zob) = \theta\chi_o/(1+\zob)$, pinned at
the target plane for every halo along the line of sight. No
normalisation is needed: the observable is additive over
neighbour halos, so the conditional expectation around one
target is already a mean per cluster. Every
projected structure that lenses the source galaxies behind the
target is weighted by exactly the $\brm(\theta)$ defined above:
an error in the selection-bias plateaus is an error in the
cluster lensing profile, and hence in the mass calibration.

The plateaus are not free functions --- the remainder of this
section shows they are fixed by requiring the projected-galaxy
budget of the sample to close.

\paragraph{Deriving the scale-dependent bias.} The observed
richness can be split as
\begin{equation}
\lob \;=\; \ltr \;+\; \Dp,
\label{eq:lob_split}
\end{equation}
with $\dprj$ the number of \emph{projected} galaxies: members of
other halos in the target's photo-$z$ cylinder whose sky positions
fall inside the target's aperture. Specifically, for clusters at
redshift $\zob$ with observed richness $\lob$ and true richness
$\ltr$, the mean contamination due to projected systems reads
\begin{equation}
\Dp(\lob,\ltr,\zob)
= 2\pi\int_0^{\pi} d\theta \sin\theta \;
\Dl_\theta \,,
\label{eq:Dprj_outer}
\end{equation}
where $\Dl_\theta$ is the mean projected richness per unit solid
angle contributed at angular separation $\theta$ from the target
centre. It encodes the \textsc{redMaPPer}--halo association over
the photo-$z$ cylinder:
\begin{align}
\Dl_\theta
\;&\equiv\; \Dlrnd
\;+\; \brmss\,[1-\sigma(\theta)]\; \tDllss{\theta}
\;+\; \brmls\,\sigma(\theta)\; \tDllss{\theta}
\label{eq:Dlam_split}\\
&=\int\!dz\,\frac{dV}{dz\,d\Omega}(z)\,
w(z,\zob) \int\!dM\,n(M,z)\,
\big[1+\xi_{\rm rm,h}(\theta;M,z)\big]
\notag\\
&\hspace{9em}\times
\int_0^{\lob}\!d\lambda\,P(\lambda\mid M, z)\,\lambda\,
f_A\big(\theta,\thob,\theta_\lambda(\lambda,z)\big),
\label{eq:Pop}
\end{align}
reading right to left: each halo of mass $M$ at redshift $z$
hosts $\lambda$ galaxies with probability $P(\lambda\mid M,z)$; a
fraction $f_A$ of them falls inside the target aperture; halos
are counted by the mass function $n(M,z)$ and placed either
uniformly along the cylinder (the $1$) or clustered around the
target (the $\xi_{\rm rm,h}$); the photo-$z$ window $w$ selects
in redshift. Two definitions:
\begin{itemize}
\item \textbf{Eligibility ($\lambda<\lob$).}
      \textsc{redMaPPer} processes halos in decreasing richness
      order and lets richer systems claim member galaxies first,
      so only halos with $\lambda<\lob$ still have galaxies left
      to contaminate the target --- the upper limit of every
      $\lambda$-integral is the percolation cut, not a
      convenience.
\item \textbf{Capture fraction $f_A$.} A projected halo's
      $\lambda$ galaxies are spread uniformly over its own disk
      of angular radius $\theta_\lambda(\lambda,z)$; the target
      aperture (radius $\thob$) captures the fraction
      $f_A = A_{\rm ov}/(\pi\theta_\lambda^2)\in[0,1]$, with
      $A_{\rm ov}$ the two-disk overlap area at centre separation
      $\theta$. $f_A$ vanishes for
      $\theta>\thob+\theta_\lambda$, so the support of
      eq.~\eqref{eq:Dprj_outer} ends at $\theta\le 2\thob$
      despite the formal $\pi$ upper limit. Because
      $\theta_\lambda$ depends on the projected halo's own
      $\lambda$, $f_A$ sits inside the $\lambda$-integral.
\end{itemize}

Substituting the correlation model
eqs.~\eqref{eq:xi_rmh}--\eqref{eq:brm} into eq.~\eqref{eq:Pop}
produces exactly the decomposition announced in
eq.~\eqref{eq:Dlam_split}, with the two channel profiles defined
by the weight they carry under the same
$(z,M,\lambda)$-integral:
\begin{align}
\Dlrnd
&= \int\!dz\,\frac{dV}{dz\,d\Omega}\,w
\int\!dM\,n
\int_0^{\lob}\!d\lambda\,P\,\lambda
&&\text{(weight $1$: uniform interlopers)},
\label{eq:Dlam_rnd}\\
\tDllss{\theta}
&= \int\!dz\,\frac{dV}{dz\,d\Omega}\,w
\int\!dM\,n\; b\,\ximm
\int_0^{\lob}\!d\lambda\,P\,\lambda\,f_A
&&\text{(weight $b\,\ximm$: clustered structures)},
\label{eq:Dlam_lss}
\end{align}
where the tilde marks that the target's own bias $\brm$ has been
divided out ($\brm$ set to one): $\tDllss{\theta}$ is a pure
integral, computable without knowing the selection bias, and the
unknown plateaus $\brmss$, $\brmls$ multiply it from outside.
Note that the uniform channel carries no capture fraction and no
$\theta$: smearing a uniform field over the donors' disks leaves
it uniform, and $f_A$ integrates out exactly ---
$\int d^2\theta\, f_A = \pi\thob^2$ for every donor size (the
total-donation identity) --- so $\Dlrnd$ is the mean eligible
richness per unit solid angle in the cylinder, constant in
$\theta$, and comes out of the sky integral of
eq.~\eqref{eq:Dprj_outer}.

\paragraph{Assembled contamination.} Inserting
eq.~\eqref{eq:Dlam_split} into eq.~\eqref{eq:Dprj_outer} and
integrating over the sky:
\begin{equation}
\Dp \;=\; \Dprnd
\;+\; \Dpss
\;+\; \Dpls,
\label{eq:Dprj_split}
\end{equation}
with the three members
\begin{align}
\Dprnd
&= \pi\,\thob^2\; \Dlrnd,
\label{eq:Dprj_rnd}\\
\Dpss
&= \brmss\; \tDpss,
\qquad
\tDpss
= 2\pi\!\int_0^\pi\! d\theta\sin\theta\;
[1-\sigma(\theta)]\; \tDllss{\theta},
\label{eq:Dprj_ss}\\
\Dpls
&= \brmls\; \tDpls,
\qquad
\tDpls
= 2\pi\!\int_0^\pi\! d\theta\sin\theta\;
\sigma(\theta)\; \tDllss{\theta}.
\label{eq:Dprj_ls}
\end{align}
$\Dprnd$ counts the uniform (random) interlopers --- the constant
$\Dlrnd$ times the aperture area, per the total-donation identity
above; $\tDpss$ and $\tDpls$ are the small- and large-scale
halves of the clustered (large-scale-structure) contamination,
computed numerically with unit selection bias. In the production
code these are \texttt{P1}, $I_2-I_1$ and $I_1$ respectively.

\paragraph{Closing the selection bias.} The plateaus follow from
two conditions. First, the large-scale plateau is anchored to the
halo bias of a mass-selected sample at the same $(\lob,\zob)$,
\begin{equation}
\bhalo(\lob,\zob)
= \frac{\int d\ln M\;M\,n(M,\zob)\,b(M,\zob)\,P(\lob\mid M,\zob)}
       {\int d\ln M\;M\,n(M,\zob)\,P(\lob\mid M,\zob)},
\label{eq:bhalo}
\end{equation}
through the Buzzard-calibrated empirical relation
\begin{align}
\brmls \;&=\; \bhalo \left[ 1 + 0.13\, \dprjx\right],
\qquad
\dprjx
= \frac{\lob-\ltr}{\Dp_{\rm halo}} - 1,
\notag\\
\Dp_{\rm halo} \;&=\; \Dprnd
+ \bhalo\big[\tDpss + \tDpls\big],
\label{eq:bls}
\end{align}
where $\Dp_{\rm halo}$ is the contamination a \emph{mass-selected}
target would suffer --- eq.~\eqref{eq:Dprj_split} with the constant
$\bhalo$ in place of both plateaus --- and $\dprjx$ measures the
fractional excess of the observed inflation over it (the
calibration of the $0.13$ slope is tied to this denominator).
Second, the observed inflation must
be exactly accounted for: substituting
eqs.~\eqref{eq:Dprj_rnd}--\eqref{eq:Dprj_ls} into
eqs.~\eqref{eq:lob_split} and \eqref{eq:Dprj_split},
\begin{equation}
\lob-\ltr = \Dprnd
+ \brmls\,\tDpls
+ \brmss\,\tDpss,
\label{eq:budget}
\end{equation}
which closes the small-scale plateau as the remaining linear
unknown:
\begin{equation}
\boxed{\;
\brmss \;=\;
\frac{(\lob-\ltr) - \Dprnd - \Dpls}
     {\tDpss}.
\;}
\label{eq:bss}
\end{equation}

\paragraph{Algorithm.} The computation at fixed
$(\lob,\ltr,\zob)$ is therefore four numerical steps and one
closure:
\begin{enumerate}
\item compute $\bhalo(\lob,\zob)$, eq.~\eqref{eq:bhalo} --- a
      1-D mass integral;
\item compute $\Dprnd$, eq.~\eqref{eq:Dprj_rnd} --- the constant
      $\Dlrnd$ times the aperture area;
\item compute the unit-bias clustered integrals $\tDpss$ and
      $\tDpls$, eqs.~\eqref{eq:Dprj_ss}--\eqref{eq:Dprj_ls};
\item form $\Dp_{\rm halo}$, $\dprjx$ and the large-scale plateau
      $\brmls$, eq.~\eqref{eq:bls};
\item close the small-scale plateau $\brmss$,
      eq.~\eqref{eq:bss}.
\end{enumerate}
With $\brmss$ and $\brmls$ in hand, eq.~\eqref{eq:brm} returns
the full scale-dependent selection bias
$\brm(\theta,\lob,\zob)$.

\section{The frozen-physics algorithm}
\label{sec:frozen_algo}

The operators of \S\ref{sec:question} are 4-D quadratures over
$(z, M, \lambda, \theta)$, and in the original variables the
integrand is numerically hostile: the halo-exclusion condition,
sliced along $z$, produces two sharp peaks bridged by a plateau
in the $z$-integrand, on a scale $\delta z\sim 10^{-3}$. This
section derives the equivalent --- but numerically friendly ---
forms actually implemented, built on one physical observation and
one change of variables.

\subsection{Freezing the slow redshift dependences}
\label{sec:freeze}

The redshift barely varies across the scales that matter. The
redMaPPer photo-$z$ window, $\sigma_z(z)$ of \texttt{photoz.py}
(the DES-Y3 richness-estimator kernel, read off a calibration
table --- not a single-galaxy photo-$z$ error), has width
$\sigma_z\simeq0.09$ at $\zob=0.5$, i.e.\ comoving
$\sigma_\chi\simeq205\,\hMpc$
($\sigma_\chi\equiv\tfrac12[\chi(\zob+\sigma_z)-\chi(\zob-\sigma_z)]$).
It is much wider than the naive single-galaxy estimate
$\sigma_z^{\rm gal}\sim0.01(1+z)\simeq0.013$ at $z=0.3$
($\sigma_\chi^{\rm gal}\simeq34\,\hMpc$): $\sigma_z(z)$ is the
scatter of the richness \emph{estimator} itself, built from many
member galaxies and their individual photo-$z$'s, not the
redshift error of any one galaxy. Meanwhile the exclusion
structure lives on
$\Rex = R_\lambda(\lob)(1+\zob)\simeq1.1\,\hMpc$
($\delta z\simeq5\times10^{-4}$); the astrophysical weights ---
$n(M,z)$, $b(M,z)$, $P(\lambda\mid M,z)$, the aperture
$\theta_\lambda(\lambda,z)$ --- evolve on $\Delta z\sim0.1$--$1$.
We therefore \emph{freeze} the slow weights at $\zob$ inside the
exclusion-affected zone ($|\chi(z)-\chi_o|\lesssim20\Rex$,
i.e.\ $\Delta z\lesssim10^{-2}$; $\chi_o\equiv\chi(\zob)$). Two
protections control the error: the drift itself is
$\partial_z\ln(\cdot)\,\Delta z\lesssim10^{-2}$, and its leading
term is \emph{odd} in $\chi(z)-\chi_o$ while the correlation
kernel is even, so the linear error integrates to zero and only
the quadratic survives, $\sim10^{-4}$. Direct validation against
\texttt{scipy.quad} puts the freeze at $0.001$--$0.003\%$ on the
clustered operators across the DES-Y3 range. The photo-$z$
window $w(z,\zob)$ is \emph{never} frozen --- it is the one
genuine line-of-sight structure and stays explicit in every
integral below.

\subsection{Sky variables}
\label{sec:sky_vars}

Both sharp structures --- the exclusion ball and the correlation
function --- depend only on the 3-D separation between target and
halo, while the aperture is transverse and the photo-$z$ window
is line-of-sight. Decompose the separation accordingly: a
line-of-sight component and an exact sky-plane (transverse)
component,
\begin{equation}
\pi \equiv \chi(z)-\chi_o,
\qquad
s \equiv 2\sqrt{\chi\chi_o}\,\sin(\theta/2)
\quad\Longrightarrow\quad
r^2 = \pi^2 + s^2 \;\;\text{(exact)},
\label{eq:sky_vars}
\end{equation}
where $r=|\bm{r}_{\rm rm}-\bm{r}_{\rm h}|$ is the comoving
separation entering $\ximm$. The measure maps exactly as well:
\begin{equation}
dz\,\frac{dV}{dz\,d\Omega}\;2\pi\sin\theta\,d\theta
= 2\pi\Big(1+\tfrac{\pi}{\chi_o}\Big)\,d\pi\;s\,ds
\label{eq:measure}
\end{equation}
(the $2\pi$ prefactors are the number; the line-of-sight variable
$\pi$ always carries a differential or an argument). In these
variables the geometry is transparent. The aperture support
$f_A=0$ beyond $\theta_\lambda+\thob\le2\thob$ becomes the
cylinder $s\le2\Rex$; the halo-exclusion condition
($\ximm\equiv0$ for separations below $\Rex$,
\S\ref{sec:question}) becomes the clean spherical statement
$r>\Rex$ --- a hard \emph{lower limit} of the radial integral,
not a discontinuity inside the domain. The twin peaks of the
original $z$-integrand were never physical structure: they are
the shadow of this excluded ball sliced along $z$, and they never
form in $(\pi,s)$.

\subsection{Reduced operators}
\label{sec:operators}

At each $(\lob,\zob)$ the algorithm evaluates the random channel
and the clustered channel, the latter split into two
line-of-sight zones:
\begin{equation}
\Dprnd = \pi\,\thob^2 \int\!dz\;\frac{dV}{dz\,d\Omega}\;
w(z,\zob)\;\Dl(z),
\qquad
\tDpss = \tDp^{\rm ss,excl}_{\rm lss} + \tDp^{\rm ss,cyl}_{\rm lss}.
\label{eq:alg_split}
\end{equation}
(The fiducial pipeline additionally removes the halo-exclusion
sphere from the random channel; in the sky variables this
carve-out is the incomplete transverse moment
$M_0(2)-M_0(x_{\rm excl}(\pi))$, with
$M_0(y)=\int_0^y x\,I_0(x)\,dx$ over the unweighted analogue of
eq.~\ref{eq:Ilam} and
$x_{\rm excl}=\sqrt{\max(0,\,1-(\pi/\Rex)^2)}$ --- a ring-sliver
correction of a few $\times10^{-3}$ relative, adopted by the
implementation so that all three operators share one exclusion
convention.) The \emph{exclusion zone}, $|\pi|\le\pis\equiv20\,\Rex$, is the
cluster's own neighbourhood: it contains the halo-exclusion
sphere $r<\Rex$ and the target aperture, and every calculation
there must track the boundary $r=\Rex$. The
\emph{free-of-exclusion zone} is the rest of the photo-$z$
cylinder, out to the edges of the window: beyond $\pis$,
$r\ge|\pi|>\pis=20\Rex\gg\Rex$ automatically, so the exclusion
condition is satisfied by every point and never has to be
imposed --- there is no boundary left to track. Both zones stay
inside the transverse cylinder $s\le2\Rex$; only the exclusion
\emph{sphere} is absent from the second. Throughout we write the
small-scale (ss) operator; the large-scale one is the same
expression with $\Ilamss\to \Ilamls$.

The split point balances two competing requirements. The
exclusion zone must stay narrow enough in redshift that the
astrophysical weights can be frozen at $\zob$
($\Delta z\lesssim10^{-2}$, \S\ref{sec:freeze}); the
free-of-exclusion zone must start far enough out that
$s\ll|\pi|$ and the transverse expansion of
eq.~\eqref{eq:taylor} converges. At $\pis=20\Rex$ both are
simultaneously controlled: the largest ratio in the
free-of-exclusion zone is $(2\Rex/\pis)^2 = 10^{-2}$, the
$\Bss$ term of eq.~\eqref{eq:cyl} retains the first
correction at exactly this order, and the next omitted
contribution is $(2\Rex/\pis)^4\lesssim10^{-4}$. At the
reference $\Rex\simeq1.1\,\hMpc$ the geometry is
$s_{\max}\simeq2.2\,\hMpc$, $\pis\simeq22\,\hMpc$,
$\rmax=\sqrt{22^2+2.2^2}\simeq22.1\,\hMpc$: the exclusion
zone is a short central segment of a much longer photo-$z$
cylinder ($\sigma_\chi\simeq205\,\hMpc$ per side).

The building block of both regimes is the frozen mass-integral
moment (Fig.~\ref{fig:profiles}),
\begin{equation}
\boxed{\;
\Ilam(x) = \int\!dM\, n(M,\zob)\,b(M,\zob)
\int_0^{\lob}\!d\lambda\, P(\lambda\mid M,\zob)\,\lambda\,
f_A(x,x_\lambda),
\;}
\quad
x = \tfrac{s}{\Rex},
\;\;
x_\lambda = \big(\tfrac{\lambda}{\lob}\big)^{0.2}\!,
\label{eq:Ilam}
\end{equation}
with the channel versions carrying the sigmoid weight,
\begin{equation}
\Ilamss(x) = [1-\sigma(x)]\,\Ilam(x),
\qquad
\Ilamls(x) = \sigma(x)\,\Ilam(x).
\label{eq:Ilam_ssls}
\end{equation}
Two contamination densities accompany it: the unweighted one that
drives the random channel, and its bias-weighted counterpart that
sets the amplitude in regime B,
\begin{align}
\Dl(z) &= \int\!dM\,n(M,z)\!\int_0^{\lob}\!d\lambda\,
P(\lambda\mid M,z)\,\lambda,
\notag\\
\Dlb(z) &= \int\!dM\,n(M,z)\,b(M,z)\!\int_0^{\lob}\!d\lambda\,
P(\lambda\mid M,z)\,\lambda.
\label{eq:densities}
\end{align}
These are not new objects. Eligibility makes every donor disk
smaller than the target aperture ($x_\lambda\le1$), so
$f_A(0,x_\lambda)=1$: at zero transverse offset the mass-integral
moment captures every donated galaxy, and
\begin{equation}
\Ilam(0) = \Dlb(\zob)
\label{eq:Ilam_zero}
\end{equation}
--- the densities are the $\theta=0$, $z=\zob$ limits of the
frozen profiles ($\Dl(\zob)$ likewise from the unweighted
analogue of eq.~\ref{eq:Ilam}).

\begin{figure}[!htbp]
\centering
\includegraphics[width=\textwidth]{figs/pedag_frozen_kernels.png}
\caption{\emph{Left:} the frozen mass-integral moment at
$(\lob,\zob)=(20,0.5)$: $\Ilam(x)$ (red, normalised to
$x=0$), its large-scale part $\Ilamls=\sigma
\Ilam$ (blue), and the sigmoid $\sigma(x)$ (dashed);
$\Ilamss$ is the difference of the two curves.
Support ends at $x=1+x_{\lambda,\max}\le2$. \emph{Right:} the
capture fraction $f_A(x,x_\lambda)$ of \S\ref{sec:question} for
several donor sizes; $\Ilam$ is their bias- and
richness-weighted average.}
\label{fig:profiles}
\end{figure}

\paragraph{Exclusion zone.} The halo-exclusion sphere is
spherical in $(\pi,s)$, so integrate on the sphere: with
$\mu=\pi/r$,
\begin{equation}
\boxed{\;
\tDp^{\rm ss,excl}_{\rm lss}
= 2\pi \int_{\Rex}^{\rmax}\! dr\; r^2\,\ximm(r)\;
\Gss(r),
\qquad
\Gss(r) = 2\!\int_{\mulo(r)}^{\muhi(r)}\! d\mu\;
\Ilamss\Big(\tfrac{r}{\Rex}\sqrt{1-\mu^2}\Big)\,
\barw(r\mu),
\;}
\label{eq:excl}
\end{equation}
where $\barw$ is the two-sided average of the photo-$z$ window
and the $\mu$-limits trace the domain boundaries: the cylinder
wall gives $\mulo(r)=\sqrt{\max(0,\,1-(2\Rex/r)^2)}$, the
zone split gives $\muhi(r)=\min(1,\,\pis/r)$, and the
radial integral ends at $\rmax=\sqrt{\pis^2+4\Rex^2}$, the
corner of the domain. The integrand is smooth
(Fig.~\ref{fig:radial_new}): the exclusion is the hard lower
limit $r=\Rex$, and the only non-analytic point is a derivative
kink at $r=2\Rex$ where the cylinder wall engages --- split the
grid there and log-spaced Gauss--Legendre nodes converge
everywhere.

\begin{figure}[!htbp]
\centering
\includegraphics[width=0.72\textwidth]{figs/pedag_frozen_radial.png}
\caption{Exclusion-zone radial integrand
$r^2\,\ximm(r)\,G(r)$ of eq.~\eqref{eq:excl} at
$(\lob,\zob)=(20,0.5)$, for the total (ss+ls) and large-scale
mass integrals. Exclusion radius = hard lower limit (dotted);
derivative kink at $r=2\Rex$ (dashed). The twin peaks and ring
plateau of the original $z$-integrand are absent.}
\label{fig:radial_new}
\end{figure}

\paragraph{Free-of-exclusion zone.} Beyond $\pis$ the geometry
is Limber-like: the cylinder is \emph{thin}. Its radius,
$2\Rex\simeq2\,\hMpc$, is small against every scale that varies
along the line of sight ($|\pi|\ge\pis\simeq22\,\hMpc$, photo-$z$
window $\sigma_\chi\simeq205\,\hMpc$), so expand the separation
about the cylinder axis,
\begin{equation}
r = |\pi|\sqrt{1+s^2/\pi^2}
= |\pi| + \frac{s^2}{2|\pi|} + O\!\big(s^4/|\pi|^3\big),
\qquad
\ximm(r) = \ximm(|\pi|)
+ \frac{s^2}{2|\pi|}\,\ximm'(|\pi|)
+ O\!\Big[\big(\tfrac{2\Rex}{\pi}\big)^{4}\Big].
\label{eq:taylor}
\end{equation}
This is the Limber step: the correlation across the thin
cross-section is replaced by its value and slope on the axis.
With the exclusion sphere gone (no $r=\Rex$ boundary to track
here), the $s$-integral of eq.~\eqref{eq:step2} acts directly on
this expansion and reduces to two fixed moments of the mass
integral,
\begin{align}
\int_0^{2\Rex}\! s\,ds\; \Ilamss(x)\,\ximm(r)
&= \Ass\,\ximm(|\pi|)
+ \frac{\Bss}{2|\pi|}\,\ximm'(|\pi|)
+ O\!\Big[\big(\tfrac{2\Rex}{\pi}\big)^{4}\Big],
\notag\\
\Ass = \Rex^2\!\int_0^2\! x\,\Ilamss(x)\,dx,
&\qquad
\Bss = \Rex^4\!\int_0^2\! x^3\,\Ilamss(x)\,dx.
\label{eq:moments}
\end{align}
Two $z$-dependences remain to be restored before this can be
integrated down the line of sight. Geometric: the aperture is
fixed in \emph{angle}, so the transverse variable dilates,
$x=(s/\Rex)\sqrt{\chi_o/\chi}$, and each moment above picks up a
power of $\chi/\chi_o$ per two powers of $s$ --- combined with
the $(1+\pi/\chi_o)$ of the measure eq.~\eqref{eq:measure}, this
is the squared distance ratio below. Astrophysical: the weights
are no longer frozen, but their \emph{shape} drifts far more
slowly than their \emph{amplitude}, so keep the amplitude exactly
and freeze only the shape,
\begin{equation}
\Ilam(x;z) \;\simeq\;
\frac{\Dlb(z)}{\Dlb(\zob)}\; \Ilam(x;\zob),
\label{eq:amp_drift}
\end{equation}
anchored by eq.~\eqref{eq:Ilam_zero}
($\Ilam(0;z)=\Dlb(z)$ exactly). Writing the background side
($z>\zob$, at $z_+ = z(\chi_o+\pi)$) and the foreground side
($z<\zob$, at $z_- = z(\chi_o-\pi)$) explicitly:
\begin{align}
\tDp^{\rm ss,cyl}_{\rm lss}
=\;& 2\pi\int_{\pis}^{\chi(z_{\rm hi})-\chi_o}\! d\pi\;
\Big(1+\tfrac{\pi}{\chi_o}\Big)^{\!2}\,
w(z_+,\zob)\;
\frac{\Dlb(z_+)}{\Dlb(\zob)}
\left[ \Ass\,\ximm(\pi)
+ \frac{\Bss}{2\pi}\,\ximm'(\pi)\right]
\notag\\
+\;& 2\pi\int_{\pis}^{\chi_o-\chi(z_{\rm lo})}\! d\pi\;
\Big(1-\tfrac{\pi}{\chi_o}\Big)^{\!2}\,
w(z_-,\zob)\;
\frac{\Dlb(z_-)}{\Dlb(\zob)}
\left[ \Ass\,\ximm(\pi)
+ \frac{\Bss}{2\pi}\,\ximm'(\pi)\right],
\label{eq:cyl}
\end{align}
where $z_{\rm lo}$ and $z_{\rm hi}$ are the edges of the
photo-$z$ window (where $w$ vanishes). The zone split at
$\pis=20\Rex$ makes the expansion parameter of
eq.~\eqref{eq:taylor} equal to $(2\Rex/\pis)^2=10^{-2}$ at the
split, decreasing outward; with the first moment correction
retained, the residual is $O(10^{-4})$. (The channel sum of the
leading moments is closed analytically by the total-donation
identity, $\Ass+\Als=\Rex^2\,\Dlb(\zob)/2$ --- a
free consistency check on the numerics.) The free-of-exclusion
zone cannot be truncated early (Fig.~\ref{fig:farzone_new}): the
BAO feature sits well inside the photo-$z$ window.

\begin{figure}[!htbp]
\centering
\includegraphics[width=0.72\textwidth]{figs/pedag_frozen_farzone.png}
\caption{Free-of-exclusion-zone line-of-sight integrand of
eq.~\eqref{eq:cyl} (background side, log-$\pi$ measure). The BAO
feature (dotted) lies inside the photo-$z$ window (dashed): the
far tail carries real weight, hence both integrals run to the
edges of the window.}
\label{fig:farzone_new}
\end{figure}

\paragraph{Derivation.} Combining
eqs.~\eqref{eq:Dlam_lss} and \eqref{eq:Dprj_ss}:
\begin{align}
\tDpss
= 2\pi\!\int_0^\pi\! d\theta\sin\theta\; [1-\sigma(\theta)]
\int\!dz\,\frac{dV}{dz\,d\Omega}\,&w(z,\zob)
\int\!dM\,n(M,z)\, b(M,z)\,\ximm(r)
\notag\\[-2pt]
&\times
\int_0^{\lob}\!d\lambda\,P(\lambda\mid M,z)\,\lambda\, f_A.
\label{eq:step1}
\end{align}
The sky variables
eqs.~\eqref{eq:sky_vars}--\eqref{eq:measure} are exact and give
\begin{equation}
\tDpss
= 2\pi\!\int\! d\pi\,\Big(1+\tfrac{\pi}{\chi_o}\Big)\,
w\big(z(\pi),\zob\big)
\int_0^{2\Rex}\! s\,ds\; [1-\sigma]\,
\ximm\big(\sqrt{\pi^2+s^2}\big)
\int\!dM\,n\,b
\int_0^{\lob}\!d\lambda\,P\,\lambda\, f_A.
\label{eq:step2}
\end{equation}

\emph{Freeze.} For $|\pi|\le\pis$ the slow weights are
evaluated at $\zob$ (\S\ref{sec:freeze}). $\ximm$ carries no
$M$-dependence, so it passes through the mass integral and the
$(M,\lambda)$ block of eq.~\eqref{eq:step2} contracts, with the
sigmoid universal in $x$:
\begin{equation}
[1-\sigma(\theta)]
\int\!dM\,n\,b
\int_0^{\lob}\!d\lambda\,P\,\lambda\, f_A
\;\;\xrightarrow{\;z\,\to\,\zob\;}\;\;
[1-\sigma(x)]\, \Ilam(x) \;=\; \Ilamss(x).
\label{eq:contract}
\end{equation}

\emph{Exclusion zone.} The frozen integrand depends on $(\pi,s)$
only through $r=\sqrt{\pi^2+s^2}$ and $x=s/\Rex$, i.e.\ through
$r$ and the direction cosine $\mu=\pi/r$. The spherical map
$d\pi\,s\,ds = r^2\,dr\,d\mu$ turns the domain
$\{r>\Rex,\;s\le2\Rex,\;|\pi|\le\pis\}$ into the limits
$\mulo(r)$, $\muhi(r)$, $\rmax$ of
eq.~\eqref{eq:excl}; the factor $(1+\pi/\chi_o)$ is odd about
the cluster and averages to one by parity, and the two
line-of-sight signs fold into $\barw$. This is
eq.~\eqref{eq:excl}.

\emph{Free-of-exclusion zone.} Here $s\le2\Rex\ll|\pi|$: Taylor-
expanding $r$ and $\ximm(r)$ about the cylinder axis
(eq.~\ref{eq:taylor}, the Limber step) turns the $s$-integral of
eq.~\eqref{eq:step2} into the fixed moments $\Ass$,
$\Bss$; restoring the geometric dilation of $x$ and the
amplitude-drift ansatz eq.~\eqref{eq:amp_drift} for the
astrophysical weights gives eq.~\eqref{eq:cyl}. See
\S\ref{sec:operators} for the expansion, the two $z$-restorations,
and the error budget in full.

\emph{Random channel.} No $\ximm$, no exclusion: $f_A$
integrates out over the sky (total-donation identity) and the
first member of eq.~\eqref{eq:alg_split} follows directly from
eqs.~\eqref{eq:Dlam_rnd} and \eqref{eq:Dprj_rnd}.

The full validation --- operators against \texttt{scipy.quad},
marginalised plateaus and the end-to-end $\brm(\theta)$ against
the production fiducial --- is \S\ref{sec:validation}.

\section{Validation}
\label{sec:validation}

Reference implementation: \texttt{FrozenSelBias}
(\texttt{src/richness\_selection/frozen\_bsel.py}), a drop-in for
the production \texttt{SelBias} --- only the computation of the
three operators changes; the closure, the $\ltr$-marginalisation
and the lensing pipeline are shared code. Both methods use the
fiducial (Costanzi) conventions --- the $\Dp_{\rm halo}$
denominator of eq.~\eqref{eq:bls} and the random-channel
exclusion carve-out of \S\ref{sec:operators} --- so every
residual below is pure operator numerics. Scripts:
\texttt{validations/frozen\_bsel\_validation.py} (operators,
plateaus, budget),
\texttt{docs/make\_validation\_plots\_frozen\_physics.py}
(end-to-end figures).

\subsection{Operators against \texorpdfstring{\texttt{scipy.quad}}{scipy.quad}}

Adaptive quadrature in $z$ ($\epsilon_{\rm rel}=10^{-6}$) on the
original 4-D form eq.~\eqref{eq:step1}, matched inner grids,
against the production grid code (fiducial) and the frozen
algorithm:
\begin{center}
\small
\begin{tabular}{r r l r r r}
\toprule
$\lob$ & $\zob$ & source & $\Dprnd$ err & $\tDpls$ err
& (ss+ls) err \\
\midrule
$20.0$  & $0.500$ & fiducial & $0.018\%$ & $0.004\%$ & $0.002\%$ \\
        &         & frozen   & $0.006\%$ & $0.037\%$ & $0.037\%$ \\
$52.5$  & $0.425$ & fiducial & $0.007\%$ & $0.020\%$ & $0.021\%$ \\
        &         & frozen   & $0.015\%$ & $0.030\%$ & $0.019\%$ \\
$130.0$ & $0.575$ & fiducial & $0.004\%$ & $0.026\%$ & $0.027\%$ \\
        &         & frozen   & $0.017\%$ & $0.003\%$ & $0.003\%$ \\
\bottomrule
\end{tabular}
\end{center}
Worst frozen error $0.037\%$, comfortably below the $0.1\%$
pipeline tolerance and on par with the fiducial grid code.

\subsection{Marginalised plateaus and budget closure}

The $\ltr$-marginalised plateaus (the scalars that feed the
lensing kernel through eq.~\ref{eq:brm}), fiducial
$\rightarrow$ frozen:
\begin{center}
\small
\begin{tabular}{r r c c r}
\toprule
$\lob$ & $\zob$ & $\brmss$ & $\brmls$
& max per-$\ltr$ dev \\
\midrule
$20.0$  & $0.500$ & $23.476 \rightarrow 23.484$
& $3.824 \rightarrow 3.824$ & $0.06\%$ \\
$52.5$  & $0.425$ & $21.056 \rightarrow 21.052$
& $5.924 \rightarrow 5.924$ & $0.03\%$ \\
$130.0$ & $0.575$ & $55.933 \rightarrow 55.957$
& $12.440 \rightarrow 12.440$ & $0.04\%$ \\
\bottomrule
\end{tabular}
\end{center}
The per-$\ltr$ plateau vectors deviate by at most $0.06\%$
anywhere on the marginalisation grid. The budget identity
eq.~\eqref{eq:budget} closes to machine precision
($\max|{\rm budget}-(\lob-\ltr)| = 1.8\times10^{-15}$) --- exact
by construction of eq.~\eqref{eq:bss}, verified as an
implementation check.

\subsection{End-to-end: \texorpdfstring{$\brm(\theta)$}{brm} and the lensing observable}

Figure~\ref{fig:bsel_vs_prod} sweeps the fixed-$\ltr$ profiles at
$(\lob,\zob)=(20,0.5)$: the frozen curves are visually coincident
with production, with residuals $\le0.044\%$ of the profile scale
--- worst for the sign-changing $\dprj=0$ profile, where the
$\brmss$ closure is a near-cancellation and amplifies operator
errors; the amplification stays mild because ss and ls share one
pipeline (\S\ref{sec:operators}) and their errors largely cancel
in eq.~\eqref{eq:bss}. Figure~\ref{fig:dsigma_vs_prod} propagates
the comparison through the lensing pipeline into
$\avg{\DSprj(R)}$ (eq.~\ref{eq:DSprj}): residuals stay below
$\simeq0.05\%$ everywhere away from the $\dprj=0$ sign crossing.

\begin{figure}[!htbp]
\centering
\includegraphics[width=0.92\textwidth]{figs/pedag_frozen_vs_prod_bsel.png}
\caption{$b_{\rm sel}(\theta\mid\lob,\ltr,\zob)$ at
$(\lob,\zob)=(20,0.5)$ for $\dprj\in\{0,3,6,10,15\}$: production
(dashed) vs frozen (solid, visually coincident). \emph{Bottom:}
residual normalised by $\max_\theta|b^{\rm prod}|$ per $\dprj$;
worst $0.044\%$ at $\dprj=0$.}
\label{fig:bsel_vs_prod}
\end{figure}

\begin{figure}[!htbp]
\centering
\includegraphics[width=0.8\textwidth]{figs/pedag_frozen_vs_prod_delta_sigma.png}
\caption{$\avg{\DSprj(R\mid\lob,\zob)}$ (cl+LSS piece) at the
same $(\lob,\zob,\dprj)$ grid, production (dashed, circles) vs
frozen (solid, squares); curves overlap at plot resolution.
\emph{Bottom:} fractional residual, $\lesssim0.05\%$ throughout.
Points within $5\%$ of the $\dprj=0$ curve's sign crossing are
dropped from the residual: the fractional ratio diverges there
while the absolute difference stays at the same level as
everywhere else.}
\label{fig:dsigma_vs_prod}
\end{figure}

\subsection{Cost}

Per $(\lob,\zob)$, median over repeats: frozen $9.2\,$ms fresh
vs production $6.4\,$ms; both $\sim0.3\,\mu$s on cached repeat
calls (all levels of the frozen computation --- the
$(M,\lambda)$ contraction, the line-of-sight tables, the drift
splines, the assembled operators --- are memoised per
$(\lob,\zob)$). Benchmark:
\texttt{validations/benchmark\_frozen.py}.

\section{Extension: frozen physics for \texorpdfstring{$\avg{\DSprj}$}{DeltaSigma-prj}}
\label{sec:frozen_dsprj}

The lensing observable eq.~\eqref{eq:DSprj} carries the same
two-tracer weight $[1 + b(M,z)\,\brm(\theta)\,\ximm(r)]$ as the
contamination operators, so the same freeze applies. Four
conventions of the implementation
(\texttt{sigma\_prj.py} / \texttt{delta\_sigma\_prj.py}) pin the
derivation:
\begin{itemize}
\item \textbf{The offset is pinned at the target plane.}
      $\Rmis(\theta) = \theta\,D_A(\zob) = \theta\chi_o/(1+\zob)$
      for every halo along the line of sight --- the kernel
      argument never dilates with the halo redshift. Since
      $\theta\chi_o = s_o(\theta)\cdot\theta/[2\sin(\theta/2)]$
      with $s_o(\theta)\equiv2\chi_o\sin(\theta/2)$ the
      transverse variable of eq.~\eqref{eq:sky_vars} evaluated at
      the target, the kernels may equivalently be axed on
      $(R,s_o)$ with $\Rmis = s_o/(1+\zob)$, up to the arc--chord
      factor $1+\theta^2/24 \le 1+5\times10^{-5}$ at
      $\theta_{\max}$.
\item \textbf{The kernel is redshift-free.} The NFW lookup uses a
      fixed concentration and today's densities
      ($r_s(M)$, $\rho_{\rm eff}(M)$ carry no $z$), so
      $\DSmis(R\mid M,\Rmis(\theta))$ depends on $z$ through
      \emph{nothing}: the mass contraction is exactly a
      $(R,\theta)$ object.
\item \textbf{The correlation argument is the exact chord.} The
      code evaluates $\ximm$ at
      $r = |\bm r_{\rm rm}-\bm r_{\rm h}|
      = \sqrt{\chi^2+\chi_o^2-2\chi\chi_o\cos\theta}$, which by
      the chord law is exactly
      $r^2 = \pi^2 + s^2$ with
      $s^2 = s_o^2(\theta)\,(1+\pi/\chi_o)$ and
      $\pi = \chi(z)-\chi_o$. It reduces to $|\pi|$ only at
      $\theta=0$; at the largest angles the lensing integral
      reaches ($s$ up to $\sim3\max R$) the transverse part
      dominates. Below we write
      $\ximm(\pi,\theta)\equiv\ximm(r)$ directly in the
      integration variables. The growth epoch is frozen at
      $\zob$ --- a production convention, not a new
      approximation.
\item \textbf{Exclusion sits on the clustered channel only.} The
      per-$z$ mask $\theta > \theta_{\rm excl}(z)$ applied to
      $\ximm$ is \emph{set-identical} to the exact-ball condition
      $r>\Rex$ (same chord law:
      $\theta>\theta_{\rm excl}(z) \Leftrightarrow
      \chi^2+\chi_o^2-2\chi\chi_o\cos\theta > \Rex^2$). The rnd
      channel carries \emph{no} exclusion here --- opposite to
      the contamination operators, where the fiducial carves the
      random channel too (\S\ref{sec:operators}). And no
      normalisation denominator appears anywhere: the observable
      is additive over neighbour halos (eq.~\ref{eq:DSprj}).
\end{itemize}

\paragraph{Reduced form.} Define the mass-contracted lensing
kernels (exact, $z$-free):
\begin{equation}
K_n(R,\theta) = \int\!dM\; n(M,\zob)\;
\DSmis\big(R\mid M,\Rmis(\theta)\big),
\qquad
K_b(R,\theta) = \int\!dM\; n(M,\zob)\,b(M,\zob)\;
\DSmis\big(R\mid M,\Rmis(\theta)\big).
\label{eq:K_lens}
\end{equation}
The rnd channel needs \emph{no freeze at all}: because the kernel
is $z$-free, the $z$-integral commutes past it exactly,
\begin{equation}
\DSprj_{\rm rnd}(R)
= 2\pi\!\int_0^{\theta_{\max}}\! d\theta\sin\theta\;
\widetilde K_n(R,\theta),
\qquad
\widetilde K_n(R,\theta) = \int\!dM\;\tilde n(M)\,
\DSmis\big(R\mid M,\Rmis(\theta)\big),
\label{eq:rnd_exact}
\end{equation}
\begin{equation}
\tilde n(M) \equiv
\int_{z_{\rm lo}}^{z_{\rm hi}}\! dz\;
\frac{dV}{dz\,d\Omega}\;
w(z,\zob)\; n(M, z)
\label{eq:ntilde}
\end{equation}
--- one 1-D contraction per $(\lob,\zob)$, computed once over
the photo-$z$ window $[z_{\rm lo}, z_{\rm hi}]$ (the production
code re-evaluates this $\theta$-independent object at every
$\theta$-node, so the reduction hoists it for free).

For the clustered channel, freeze the shape and keep the
amplitude --- the lensing analogue of eq.~\eqref{eq:amp_drift},
\begin{equation}
n(M,z)\,b(M,z) \;=\; a_b(z)\;
n(M,\zob)\,b(M,\zob)\,\big[1+\varepsilon_b(M,z)\big],
\qquad
a_b(z) =
\frac{\displaystyle\int dM\; n(M,z)\,b(M,z)\; r_s(M)}
     {\displaystyle\int dM\; n(M,\zob)\,b(M,\zob)\; r_s(M)}.
\label{eq:nb_drift}
\end{equation}
The anchor weight is $r_s(M)\propto M^{1/3}$ --- the $R_{200}$
scale --- because that is exactly how the kernel weights mass:
$\DSmis \propto r_s\,\rho_{\rm eff}$ at fixed $R/r_s$, with
$\rho_{\rm eff}$ mass-independent, so $a_b$ tracks the drift of
the kernel-weighted bias density. By construction
$\varepsilon_b(M,\zob)=0$ and
$\int dM\,n\,b\,r_s\,\varepsilon_b = 0$ at every $z$ --- unlike
eq.~\eqref{eq:amp_drift} there is no exact $f_A(0)=1$ anchor
here, hence the explicit weight. Dropping $\varepsilon_b$, the
$(z,M)$ block at fixed $\theta$ factorises into $K_b$ times the
line-of-sight factor
\begin{equation}
\boxed{\;
\Psi(\theta) =
\int_{z_{\rm lo}}^{z_{\rm hi}}\! dz\;
\frac{dV}{dz\,d\Omega}\;
w(z,\zob)\; a_b(z)\;
\ximm(\pi,\theta)\;
\Theta\big(r > \Rex\big),
\;}
\qquad
r^2 = \pi^2 + s_o^2(\theta)\Big(1+\tfrac{\pi}{\chi_o}\Big),
\label{eq:Psi_lens}
\end{equation}
where everything except $a_b$ is exact: the measure, the window,
the chord separation (the $s$--$\pi$ coupling inside $r$ is kept
--- freezing $s\to s_o$ would cost up to $10^{-2}$ on $\ximm$ at
large $\theta$ and saves nothing, since $\Psi$ is 1-D per
$\theta$ regardless), and the exclusion ball. No $\pis$ zone
split and no Limber expansion appear anywhere: the payload is
purely transverse, so the line-of-sight integral is already 1-D,
and inside $|\pi|\lesssim\pis$ the drift ansatz degenerates into
the parity-protected hard freeze of \S\ref{sec:freeze}
($a_b\to1$). $\Psi$ is smooth except for one derivative kink
where the ball detaches, at
$\theta = \arcsin(\Rex/\chi_o) \simeq \theta_{\rm excl,o}$ ---
already a breakpoint of the production $\theta$-grid.

\paragraph{Assembly.}
\begin{equation}
\boxed{\;
\avg{\DSprj(R\mid\lob,\zob)}_{\rm tot}
= 2\pi\!\int_0^{\theta_{\max}}\! d\theta\,\sin\theta\;
\Big[\, \widetilde K_n(R,\theta)
\;+\; \brm(\theta)\,\Psi(\theta)\,K_b(R,\theta) \,\Big],
\;}
\label{eq:DSprj_frozen}
\end{equation}
no denominator, with $\brm(\theta)$ the marginalised profile of
eq.~\eqref{eq:brm} multiplying the clustered term only --- the
first / second term are the production rnd / cl+LSS pieces, and
the pipeline's default return is the second. The $\theta$-grid
keeps the production breakpoint rule
($\{\theta_{\rm excl,o},\,\thob,\,2\thob,\,
\theta_R(R_i)\,\forall i,\,\theta_{\max}\}$, with
$\theta_{\max}=3\max(R)(1+\zob)/\chi_o$ adaptive for $\DSprj$):
$K_b(R,\theta)$ retains the load-bearing $\Rmis\simeq R$ kernel
peak at each $\theta_R$, $\Psi$ kinks at $\theta_{\rm excl,o}$,
and $\brm$ inflects at $\thob/2$. (In this convention
$\theta_{\rm excl,o}=\thob$ identically, so the set has one
fewer distinct entry than its notation suggests.)

\paragraph{What this buys.} In the production code the $(z,M)$
inner block at each $\theta$-node is a genuine 2-D quadrature,
because $n(M,z)\,b(M,z)$ couples the axes. The reduction
separates it: per $\theta$-node, one scalar $\Psi(\theta)$ (1-D
in $\pi$, sharing the line-of-sight tables and drift splines
already cached by \texttt{FrozenSelBias}) and batched
$(R,\theta)$ table contractions $\widetilde K_n$, $K_b$ computed
once --- the inner cost per node drops from $N_z\times N_M$ to
$N_z + N_M$, and the rnd channel is exact rather than
re-quadratured per $\theta$.

\paragraph{Error budget.}
\begin{center}
\small
\begin{tabular}{l l l}
\toprule
source & order & size \\
\midrule
kernel $z$-dependence ($\Rmis$ pinned; $r_s,\rho_{\rm eff}$
$z$-free) & exact & $0$ \\
exclusion (per-$z$ slab $\equiv$ ball, cl only) & exact (set
identity) & $0$ \\
rnd channel via $\tilde n(M)$ & exact & $0$ \\
chord $r(\pi,\theta)$ kept exact & --- & $0$ \\
$nb$ freeze, $|\pi|\lesssim\pis$ & parity-protected,
$(\partial_z\ln nb\,\Delta z)^2$ & $\sim10^{-4}$ \\
drift-shape residual $\varepsilon_b$, full window &
${\rm Cov}_M[\varepsilon_b,\DSmis]$; second order & few
$\times10^{-4}$--$10^{-3}$ \\
arc--chord (only if re-axed on $s_o$) & $\theta^2/24$ &
$\le5\times10^{-5}$ \\
quadrature ($\theta$ breakpoints, $\Psi$'s $\pi$-grid) &
inherits production tolerances & unchanged \\
\bottomrule
\end{tabular}
\end{center}
\paragraph{Implementation and measured performance.}
\texttt{FrozenDeltaSigmaPrj}
(\texttt{src/richness\_selection/frozen\_delta\_sigma\_prj.py})
subclasses the production \texttt{DeltaSigmaPrj}, inheriting the
$\theta$-grid, the $(z,M)$ context, and the marginalised
$\brm(\theta)$ --- only the inner assembly changes, so measured
residuals are the reduction alone. Two further wins come for
free once the $\theta$-loop carries no $(z,M)$ work: $\Psi$ is
evaluated on the whole $\theta$-grid in one vectorised pass, and
because $\ln\Rmis(\theta)$ is monotonic in $\theta$ at fixed
$M$, the NFW lookup collapses to \emph{one} tensor-grid spline
call per mass node --- $N_M$ calls total, versus
$N_M\times N_\theta$ in production. Measured at
$(\lob,\zob)=(20,0.5)$, $R\in[0.1,30]$ (24 points): the rnd
channel is machine-identical (exact hoist), the cl channel
agrees to $7.5\times10^{-5}$ (the drift-shape residual, inside
its a priori budget above; stable at $7$--$8\times10^{-5}$
across $\dprj\in\{0,6,15\}$), and the total to
$1.1\times10^{-5}$; wall-clock $154\,$ms $\to$ $38\,$ms per
call, a $4.1\times$ speed-up.

% ==================================================================
% EVERYTHING BELOW: previous draft, commented out during the
% section-by-section re-write. Re-enable pieces as they are revised.
% ==================================================================
\iffalse
\section{OLD DRAFT: The physical question}
\label{sec:question_old}

\textsc{redMaPPer} estimates a cluster's richness by counting
member galaxies inside a disk of angular radius $\theta_\lambda$
and a photo-$z$ window around the cluster. Any galaxy that
satisfies both cuts is counted, whether or not it belongs to the
cluster. The observed richness therefore splits as
\begin{equation}
\lob \;=\; \ltr \;+\; \dprj(\lob,\ltr,\zob),
\end{equation}
with $\dprj$ the number of \emph{projected} galaxies: members of
other halos in the target's photo-$z$ cylinder whose sky positions
fall inside the target's aperture. Following Costanzi 2026
(\texttt{papers/Costanzi2026.tex}, \S``Projection kernel'',
Fig.~1), the $i$-th projected halo donates
\begin{equation}
\boxed{\;
\rhoprj_i \;=\; \lambda_i\,w(z_i,\zob)\,f_{A,i}
\;}
\label{eq:rho_prj_i}
\end{equation}
galaxies: its own richness $\lambda_i$, times the probability
$w(z_i,\zob)$ that a galaxy at redshift $z_i$ passes the target's
photo-$z$ cut (the parabolic photo-$z$ kernel --- it controls how
far down the line of sight donors matter, and is \emph{never}
approximated in this note), times the fraction $f_{A,i}$ of its
galaxies that land inside the target's aperture (the geometric
capture fraction derived in \S\ref{sec:fA}).

\paragraph{Eligibility: the percolation cut is physics, not
notation.} Not every halo can donate. \textsc{redMaPPer} builds
its catalogue by processing halos in \emph{decreasing richness
order} and letting richer systems claim member galaxies first. A
projected structure can only push galaxies into the target's
aperture if it has not already been used up by something richer
--- i.e.\ if its own true richness satisfies $\lambda<\lob$. A
halo with $\lambda\ge\lob$ cannot contaminate a lower-ranked
target: it would already have claimed those galaxies for itself
(or a still-richer neighbour would have). Percolation also
removes donors from the target's immediate 3-D neighbourhood ---
halo exclusion, implemented as $\ximm\equiv0$ for 3-D separations
below $\Rex=R_\lambda(\lob)(1+\zob)$, with
$R_\lambda(\lambda)=(\lambda/100)^{0.2}\,\hMpc$. Throughout this
note, every $\lambda$-integral carries the upper limit $\lob$;
that limit \emph{is} the eligibility rule.

\paragraph{Everything is a two-tracer correlation.} The eligible
galaxies arrive through two physically distinct channels, and the
clean way to see both is to remember that \emph{the target
cluster is itself a biased tracer} of the matter field, with bias
$\btar(\lob,\zob)$, exactly as every donor halo is a biased
tracer with bias $b(M,z)$. A correlation between two tracers
always carries both biases,
\begin{equation}
\xi_{\rm target,donor}(r) \;=\; \btar\; b(M,z)\;\ximm(r),
\label{eq:two_tracer}
\end{equation}
with $\ximm$ the nonlinear \emph{matter--matter} correlation
(halofit). It is convenient to name the donor's half of this
product --- the standard halo-model object
\begin{equation}
\xihm(r\,;M,z) \;\equiv\; b(M,z)\,\ximm(r),
\label{eq:xihm}
\end{equation}
the donor's halo--matter cross-correlation. The two channels are
then:
\begin{itemize}
\item \textbf{Poisson channel}: donors placed at random along the
      cylinder, uncorrelated with the target --- the $b=0$ limit
      of eq.~\eqref{eq:two_tracer}. Their count is a plain volume
      integral, weight $1$.
\item \textbf{Large-scale-structure (LSS) channel} (``cylinder
      galaxies''): donors clustered around the target. Their
      \emph{excess} count over Poisson is weighted by the full
      two-tracer correlation
      $\btar\,\xihm(r;M,z)$.
\end{itemize}

\paragraph{The master count.} Both channels are instances of one
object, a weighted sum over the contamination halos --- everything
in this note is a sum or a weighted average, nothing more:
\begin{align}
\Pop[X](\lob,\zob)
=\;&
\int\!dz\,\frac{dV}{dz\,d\Omega}(z)\!
\int\!dM\,n(M,z)\!
\int_0^{\lob}\!d\lambda\,P(\lambda\mid M,z)\,\lambda\,w_z(z,\zob)
\notag\\
&\quad\cdot\;2\pi\!\int\!d\theta\,\sin\theta\;
f_A\big(\theta,\thob,\theta_\lambda(\lambda,z)\big)\,
X(M,z,\theta\mid\zob),
\label{eq:Pop}
\end{align}
reading right-to-left: capture fraction $\times$ donated richness
$\times$ MOR $\times$ halo abundance $\times$ photo-$z$ window,
summed over the cylinder. The extra weight $X$ selects the
channel:
\begin{equation}
\Nbar \equiv \Pop[1],
\qquad
\Klss \equiv \Pop[\xihm],
\qquad
\KlssInf \equiv \Pop[\xihm\,\sigma(\theta)].
\label{eq:count_defs}
\end{equation}
$\Nbar$ is the \textbf{Poisson contamination}: the expected
eligible-galaxy count if donors were unclustered. $\Klss$ is the
\textbf{LSS contamination kernel}: the count weighted by
each donor's own $\xihm$ --- deliberately \emph{one bias short}
of a physical two-tracer term, because the target's bias is the
unknown the pipeline solves for; $\btar\,\Klss$ completes the
correlation. $\KlssInf$ is $\Klss$ weighted by the sigmoid
$\sigma(\theta)=[1+e^{-k(\theta-\theta_0)}]^{-1}$ of the
$\bsel(\theta)$ ansatz (fixed constants $k=2.5/\thob$,
$\theta_0=\thob/2$, $\thob=\Rex/\chi_o$, $\chi_o=\chi(\zob)$);
since $\sigma\to1$ at large $\theta$, $\KlssInf$ is the
\emph{large-$\theta$ (LSS-asymptotic) piece} of $\Klss$ --- the
part of the clustered count that is eventually multiplied by the
large-scale plateau $\blarge$, while the complement
$\Klss-\KlssInf$ carries the small-scale plateau $\bsmall$
(\S\ref{sec:closure}). The $\theta$-integral runs over
$(0,\,2\thob]$; \S\ref{sec:moments} shows this cut is exact, not
a truncation.

This is the same Poisson $+$ clustering split the main note
already uses one level up, where the projected profile is
$\Sigma = \Sigma^{\rm 1h} + \Sigma^{\rm prj}$ --- here it is
denominated in galaxy counts instead of surface density.

\paragraph{Dictionary.} The production document and the code
predate this naming; the map is:
\begin{center}
\small
\begin{tabular}{l l l}
\toprule
this note & production doc / code & meaning \\
\midrule
$\ximm$ & $\xi_{\rm NL}$ / \texttt{XiNL} & nonlinear matter--matter correlation \\
$\xihm = b\,\ximm$ & $b\,\xi_{\rm NL}$ & donor halo--matter cross-correlation \\
$\Nbar$ & $\Pop[1]$ / \texttt{P1} & Poisson (unclustered) contamination count \\
$\Klss$ & $I_2$ / \texttt{I2} & LSS contamination kernel \\
$\KlssInf$ & $I_1$ / \texttt{I1} & large-$\theta$ (LSS) piece of $\Klss$ \\
$\btar$ & $b_{\rm halo}$ / \texttt{b\_eff} & target's own bias (mass-selected aggregate) \\
$\Nlss = \btar\Klss$ & $b_{\rm eff}\,I_2$ & LSS-clustered contamination \\
$\Nrnd = \Nbar+\Nlss$ & $\dprj_{\rm RND}$ & contamination of a random mass-matched target \\
\bottomrule
\end{tabular}
\end{center}

\paragraph{Scales.} Three scales control everything below (values
at the reference point $\lob=20$, $\zob=0.5$):
\begin{center}
\small
\begin{tabular}{l l l}
\toprule
scale & value & carrier \\
\midrule
exclusion & $\Rex \approx 1.1\,\hMpc$
  \;($\delta z_{\rm excl}\approx 4.7\times10^{-4}$)
  & $\ximm$ cutoff, aperture $f_A$, sigmoid $\sigma$ \\
photo-$z$ window & $\sigma_z \approx 0.09$
  \;($\sigma_\chi \approx 270\,\hMpc$) & $w_z$ \\
weight evolution & $\Delta z \sim 0.1$--$1$
  & $n$, $b$, $P(\lambda|M)$, $dV/dz\,d\Omega$ \\
\bottomrule
\end{tabular}
\end{center}
The ratio of the first to the third is $\sim 10^{-3}$: everything
sharp is geometric, everything astrophysical is slow. The
reformulation makes this separation explicit.

\section{Counting the eligible galaxies}
\label{sec:kernels}

\subsection{The capture fraction \texorpdfstring{$f_A$}{fA}}
\label{sec:fA}

$f_A$ carries the entire $\theta$-shape of the donation. The
model treats a donor of true richness $\lambda$ as $\lambda$
member galaxies spread \emph{uniformly} over the donor's own
richness disk, of angular radius
$\theta_\lambda \equiv \theta_\lambda(\lambda,z)$. When the donor
sits at centre separation $\theta$ from the target, the target's
aperture (radius $\thob$) captures each of those galaxies with
probability equal to the fractional overlap of the two disks ---
so the expected captured count is $\lambda f_A$ with
\begin{equation}
\boxed{\;
f_A(\theta,\thob,\theta_\lambda)
= \frac{A_{\rm ov}(\theta,\thob,\theta_\lambda)}
       {\pi\theta_\lambda^2},
\;}
\label{eq:fA_def}
\end{equation}
$A_{\rm ov}$ the area of the two disks' intersection at centre
distance $\theta$, closed-form (the standard circle--circle lens):
\begin{align}
A_{\rm ov}(\theta,\thob,\theta_\lambda)
=\;& \thob^2
\arccos\!\Big(\frac{\theta^2+\thob^2
                     -\theta_\lambda^2}{2\theta\,\thob}\Big)
+ \theta_\lambda^2
\arccos\!\Big(\frac{\theta^2+\theta_\lambda^2
                     -\thob^2}{2\theta\,\theta_\lambda}\Big)
\notag\\
&- \tfrac12\sqrt{\big[(\thob+\theta_\lambda)^2-\theta^2\big]
                 \big[\theta^2-(\thob-\theta_\lambda)^2\big]}
\label{eq:lens}
\end{align}
for $|\thob-\theta_\lambda|<\theta<\thob+\theta_\lambda$, with
$A_{\rm ov}=\pi\min(\thob,\theta_\lambda)^2$ for
$\theta\le|\thob-\theta_\lambda|$ and $0$ for
$\theta\ge\thob+\theta_\lambda$.

Two points the uniform-disk picture makes automatic:
\begin{itemize}
\item \textbf{Normalisation is by the donor's own area.} $f_A$
      is the fraction \emph{of the donor's galaxies} landing in
      the target aperture, hence the $\pi\theta_\lambda^2$
      denominator and $f_A\in[0,1]$: $f_A=1$ when the donor disk
      is swallowed whole ($\theta\le\thob-\theta_\lambda$,
      reachable because eligibility $\lambda<\lob$ makes the
      donor the smaller disk near $\zob$), and $f_A=0$ once the
      disks separate ($\theta\ge\thob+\theta_\lambda$).
      Normalising by the \emph{target} area $\pi\thob^2$ instead
      would not be a capture fraction, and its zeroth moment
      would retain the donor richness
      (cf.\ \S\ref{sec:moments}) --- the $\lambda$--$\theta$
      decoupling below would be lost.
      \texttt{richness\_selection.geometry.area\_overlap}
      implements eq.~\eqref{eq:fA_def} by always dividing by the
      \emph{second} disk's area regardless of which disk is
      larger; the appendix-A formula of Costanzi 2026 instead
      normalises by $\pi\min(\thob,\theta_\lambda)^2$, so the two
      disagree in the tail where a low-$z$ foreground donor's own
      angular disk is larger than the target's
      ($\theta_\lambda>\thob$, possible even for $\lambda<\lob$
      since angular size grows as $z\to0$). This note follows the
      code's convention throughout; flagging the mismatch rather
      than silently picking one.
\item \textbf{Homogeneity $\Rightarrow$ universality.}
      $A_{\rm ov}$ is homogeneous of degree two in its three
      angular arguments, so $f_A$ depends only on the ratio
      $\theta_\lambda/\thob$. With
      $x\equiv\theta/\thob = \rp/\Rex$ (\S\ref{sec:coords}) and
      the donor-size ratio
      $x_\lambda \equiv \theta_\lambda/\thob$
      ($=(\lambda/\lob)^{0.2}\in(0,1]$ at frozen $\zob$), it is
      one universal two-argument function $f_A(x,x_\lambda)$: a
      plateau $f_A=1$ on $x\le 1-x_\lambda$, the lens formula in
      between, and zero beyond $x=1+x_\lambda\le 2$
      (Fig.~\ref{fig:kernels}, right). The sigmoid constants also
      scale with $\thob$, so $\sigma(x)=[1+e^{-2.5(x-1/2)}]^{-1}$
      is universal too: nothing in the angular sector depends on
      $(\lob,\zob)$ individually.
\end{itemize}

\subsection{From the halo population to counting kernels}
\label{sec:Nkernels}

The counts of \S\ref{sec:question} run over halo mass and donor
richness, but neither variable survives in the answer: what
matters at each transverse offset $x$ is simply \emph{how many
galaxies the contamination halos donate there}. Because the
donor bias factorises out of the matter correlation
(eq.~\ref{eq:xihm}: $\xihm = b(M,z)\,\ximm(r)$, with $b$ carrying
all the $M$-dependence and $\ximm$ none), the $(M,\lambda)$
sector contracts once per $(\lob,\zob)$ into the
\textbf{eligible-galaxy counting kernels}
\begin{equation}
\boxed{\;
\begin{aligned}
\Ncal_{\rm rnd}(x) &\equiv \int\!dM\, n(M,\zob)
\int_0^{\lob}\!d\lambda\, P(\lambda\mid M,\zob)\,\lambda\,
f_A(x,x_\lambda),
\\
\Ncal_{\rm lss}(x) &\equiv \int\!dM\, n(M,\zob)\,b(M,\zob)
\int_0^{\lob}\!d\lambda\, P(\lambda\mid M,\zob)\,\lambda\,
f_A(x,x_\lambda),
\end{aligned}
\;}
\label{eq:N_def}
\end{equation}
leaving $\ximm(r)$ outside as a pure function of separation.
$\Ncal_{\rm rnd}(x)$ answers: \emph{per unit comoving volume of eligible
halos placed at scaled transverse offset $x$, how many of their
galaxies land inside the target aperture?} It is a
galaxy-donation profile --- the direct, population-averaged
version of Fig.~\ref{fig:kernels}: maximal at $x=0$ (donor disks
fully overlapping the aperture), decaying to zero at $x=2$ where
even the largest eligible donor loses contact. $\Ncal_{\rm lss}(x)$ is
the same count with each donor weighted by its clustering
strength $b(M,\zob)$ --- the version the LSS channel needs,
because a donor only produces \emph{excess} galaxies insofar as
it traces the target's density field. In this sense $\Ncal_{\rm lss}$ is
the workhorse of the whole selection-bias calculation: one 1-D
curve holding the entire population of contamination halos, its
richness function, and its clustering weight.

The LSS-asymptotic kernel needs no third contraction. The sigmoid
carries no $M$ or $\lambda$ dependence, so the $\KlssInf$ kernel
is \emph{exactly}
\begin{equation}
\sigma(x)\,\Ncal_{\rm lss}(x),
\end{equation}
zero extra $(M,\lambda)$ work. Figure~\ref{fig:kernels} shows
both kernels at the reference point.

\begin{figure}[!htbp]
\centering
\includegraphics[width=\textwidth]{figs/pedag_frozen_kernels.png}
\caption{\emph{Left:} the eligible-galaxy counting kernels at
$(\lob,\zob)=(20,0.5)$. $\Ncal_{\rm lss}(x)$ (red) is the bias-weighted
donation profile of eq.~\eqref{eq:N_def}; $\sigma\Ncal_{\rm lss}$ (blue)
follows for free; the sigmoid $\sigma(x)$ (dashed) is fully
analytic. Support ends at $x = 1+x_{\lambda,\max}\le 2$ --- the
$\theta_{\max}=2\thob$ cut of the original pipeline loses
nothing. \emph{Right:} the universal capture fraction
$f_A(x,x_\lambda)$ of \S\ref{sec:fA} for several donor sizes
$x_\lambda=\theta_\lambda/\thob$: unit plateau for
$x\le 1-x_\lambda$ (donor disk fully inside the target aperture),
lens-formula falloff, zero beyond $x = 1+x_\lambda$. $\Ncal_{\rm lss}$ is
the bias-and-richness-weighted average of these curves over
$x_\lambda(\lambda)$.}
\label{fig:kernels}
\end{figure}

\subsection{The total-donation theorem}
\label{sec:moments}

The uniform-disk picture yields the moments of $f_A$ without ever
touching the lens formula. Let $D_{\rm ob}$, $D_\lambda$ be the
target and donor disks (radii $\thob$, $\theta_\lambda$, areas
$A_{\rm ob}=\pi\thob^2$, $A_\lambda=\pi\theta_\lambda^2$). Write
the overlap area as a double indicator integral and use Fubini:
\begin{equation}
\int d^2\theta\; A_{\rm ov}(\theta)
= \int d^2\theta \int d^2u\;
\mathbf{1}_{\{u\in D_{\rm ob}\}}\,\mathbf{1}_{\{u-\theta\in D_\lambda\}}
= A_{\rm ob}A_\lambda,
\end{equation}
i.e.\ summed over all possible centre offsets $\theta$, every
point of the donor disk sweeps across every point of the target
disk exactly once. For the second moment, note
$\theta=\bm{u}-\bm{v}$ with $\bm{u}$ uniform on $D_{\rm ob}$ and
$\bm{v}$ uniform on $D_\lambda$ (independent, concentric), so
\begin{equation}
\int d^2\theta\; \theta^2\, A_{\rm ov}(\theta)
= A_{\rm ob}A_\lambda\,\avg{|\bm{u}-\bm{v}|^2}
= A_{\rm ob}A_\lambda\big(\avg{u^2}+\avg{v^2}\big)
= A_{\rm ob}A_\lambda\,\frac{\thob^2+\theta_\lambda^2}{2},
\end{equation}
using $\avg{u^2}=\thob^2/2$ for a uniform disk and the vanishing
cross term. With the capture-fraction normalisation
$f_A = A_{\rm ov}/(\pi\theta_\lambda^2)$ of eq.~\eqref{eq:fA_def}
and $x=\theta/\thob$, $x_\lambda=\theta_\lambda/\thob$:
\begin{equation}
\boxed{\;
\int_0^\infty 2x\,f_A(x,x_\lambda)\,dx = 1
\quad\text{for every } x_\lambda,
\qquad
\int_0^\infty 2x^3 f_A(x,x_\lambda)\,dx = \frac{1+x_\lambda^2}{2}.
\;}
\label{eq:moments}
\end{equation}
The zeroth moment is the \textbf{total-donation theorem}:
integrated over all sky positions, a donor gives
$\lambda\times(\text{target area})$ galaxy--area regardless of how
large its own disk is --- smearing the $\lambda$ galaxies over a
wider disk redistributes \emph{where} in $\theta$ the
contamination lands but cannot change its total. (Under the
target-area normalisation this total would instead read
$x_\lambda^2$, tying it to the donor size --- the sanity check
promised in \S\ref{sec:fA}.) The donation total is therefore
\emph{independent of the donor richness}: wherever the
$\theta$-integral runs over the full $f_A$ support, the
$\lambda$--$\theta$ coupling drops out entirely, leaving the
scalar aggregates
\begin{align}
\nrnd(z) &= \int\!dM\,n(M,z)\!\int_0^{\lob}\!d\lambda\,P(\lambda\mid M,z)\,\lambda
= \int_0^\infty\! 2x\,\Ncal_{\rm rnd}(x)\,dx\Big|_z,
\notag\\
\nlss(z) &= \int\!dM\,n(M,z)\,b(M,z)\!\int_0^{\lob}\!d\lambda\,P(\lambda\mid M,z)\,\lambda
= \int_0^\infty\! 2x\,\Ncal_{\rm lss}(x)\,dx\Big|_z.
\label{eq:blam}
\end{align}
The identity $\int 2x\,\Ncal_{\rm lss}\,dx = \nlss$ is verified
numerically to $0.0000\%$ (it is exact; the check validates the
implementation). The theorem also certifies the
$\theta_{\max}=2\thob$ cut: the $f_A$ support ends at
$x=1+x_\lambda\le2$, so the cut discards nothing. The same
mechanism collapses the Poisson count $\Nbar$ completely
(\S\ref{sec:P1}).

\section{Sky-plane and line-of-sight coordinates: \texorpdfstring{$(\rp,\pi)$}{(s, pi)}}
\label{sec:coords}

The eligibility physics of \S\ref{sec:question} lives on 3-D
separation: the exclusion ball ($r<\Rex$) and the clustering
kernel $\ximm(r)$ are spherically symmetric, while the aperture
$f_A$ is transverse and the photo-$z$ window $w_z$ is
line-of-sight. The natural variables are therefore a
line-of-sight separation
\begin{equation}
\pi \;\equiv\; \chi(z)-\chi_o
\end{equation}
and a sky-plane (transverse) separation $\rp$. The exact 3-D
comoving separation obeys
\begin{equation}
r^2
= \chi^2 + \chi_o^2 - 2\chi\chi_o\cos\theta
= \pi^2 + 4\chi\chi_o\sin^2(\theta/2),
\label{eq:exact_sep}
\end{equation}
so defining
\begin{equation}
\boxed{\;
\rp \;\equiv\; 2\sqrt{\chi\chi_o}\,\sin(\theta/2)
\quad\Longrightarrow\quad
r^2 = \pi^2 + \rp^2 \;\;\text{(exact)},
\;}
\label{eq:rp_def}
\end{equation}
the sky measure and the volume element also map exactly:
\begin{equation}
2\pi_{\vphantom{o}}\sin\theta\,d\theta
= \frac{2\pi_{\vphantom{o}}\,\rp\,d\rp}{\chi\chi_o},
\qquad
dz\,\frac{dV}{dz\,d\Omega} = \chi^2\,d\pi,
\qquad
\Longrightarrow\quad
dz\,\frac{dV}{dz\,d\Omega}\;2\pi\sin\theta\,d\theta
= 2\pi\Big(1+\tfrac{\pi}{\chi_o}\Big)\,d\pi\,\rp\,d\rp.
\label{eq:measure}
\end{equation}
(Here and below the $2\pi$ prefactors are the number
$\pi=3.14159\ldots$; the line-of-sight variable $\pi$ always
appears with an argument or a differential.)
Both identities in
eqs.~\eqref{eq:rp_def}--\eqref{eq:measure} carry \emph{no} error
term: the question of a ``higher-order $\cos\theta$ expansion''
dissolves, because the exact transverse variable exists. (Had one
used the naive $\rp\simeq\chi_o\theta$, the error would be
$\theta^2/24 \lesssim (2\thob)^2/24 \sim 10^{-7}$ --- also
irrelevant, but the exact variable is free.) The comoving-volume
Jacobian $\chi^2$ cancels against the sky-measure $1/(\chi\chi_o)$
up to the factor $(1+\pi/\chi_o)$, which is \emph{odd} in $\pi$
and cancels at first order under the near-zone integral
(\S\ref{sec:frozen}).

In these variables the counting domain is transparent: the
$\theta$-range maps to $\rp\in(0,2\Rex]$ (the $f_A$ support, in
transverse units), so the clustered kernels $\Klss$, $\KlssInf$ run
over an infinite \emph{cylinder} of radius $2\Rex$ along the line
of sight --- the ``cylinder galaxies'' of \S\ref{sec:question}
--- minus the excluded ball, which is now the clean spherical
condition $r=\sqrt{\pi^2+\rp^2}>\Rex$. (In the original
$(z,\theta)$ coordinates this same ball, sliced along $z$,
produced twin sharp peaks bridged by a ring plateau in the
$z$-integrand --- the numerical hazard the production ring+outer
Gauss--Legendre split was built to tame. \S\ref{sec:near} shows
the peaks never form in $(\rp,\pi)$.)

\section{Frozen physics: what may be evaluated at \texorpdfstring{$\zob$}{zob}}
\label{sec:frozen}

Inside the exclusion-affected zone $|\pi|\lesssim \pi_s = 20\Rex
\approx 22\,\hMpc$ ($\Delta z\lesssim 8\times 10^{-3}$), freeze
at $\zob$: the halo mass function $n(M,z)$, the halo bias
$b(M,z)$, the MOR $P(\lambda\mid M,z)$, the aperture angle
$\theta_\lambda(\lambda,z)$, and the measure factor
$(1+\pi/\chi_o)$. Two independent protections control the error:
\begin{enumerate}
\item \textbf{Scale separation}: each frozen weight drifts by
      $\partial_z\ln(\cdot)\times\Delta z \lesssim 10^{-2}$ across
      the zone.
\item \textbf{Parity}: the drift is linear in $\pi$ at leading
      order, while $\ximm(r)$ and the counting kernels are even in
      $\pi$; first-order errors integrate to zero over the
      symmetric zone, leaving only the quadratic term
      $\sim(\Delta z\,\partial_z\ln)^2\lesssim 10^{-4}$.
\end{enumerate}

Direct validation (freeze inside the exclusion ring, exact
elsewhere; error vs.\ \texttt{scipy.quad}, flat across ring node
counts $N=7/15/25$ --- i.e.\ a property of the approximation, not
of the quadrature):
\begin{center}
\small
\begin{tabular}{r r r r}
\toprule
$\lob$ & $\zob$ & $\KlssInf$ err & $\Klss$ err \\
\midrule
$25.0$  & $0.275$ & $0.001\%$ & $0.002\%$ \\
$52.5$  & $0.425$ & $0.001\%$ & $0.001\%$ \\
$130.0$ & $0.575$ & $0.002\%$ & $0.003\%$ \\
\bottomrule
\end{tabular}
\end{center}
The photo-$z$ kernel $w_z$ is \emph{not} frozen anywhere: it
enters every count below as an explicit factor at the
line-of-sight level --- multiplying the counting kernels
identically in $\Klss$ and $\KlssInf$ (they share one geometric
pipeline) and appearing as a bare 1-D factor in $\Nbar$. The
counting kernels themselves cannot carry $w_z$: they are
contractions over $(M,\lambda)$ at the single redshift $\zob$,
while $w_z$ depends on the donor redshift $z$ --- a different
integration variable, one level further out.

\section{The three counts}
\label{sec:counts}

\subsection{The Poisson count \texorpdfstring{$\Nbar$}{Nbar}: a 1-D integral}
\label{sec:P1}

The Poisson channel carries no $\ximm$ and no exclusion, so the
total-donation theorem applies at every $z$ without any zone
split: the $\theta$-integral of $f_A$ is exactly the target's
aperture area, for every donor. The 4-D count collapses to
\begin{equation}
\boxed{\;
\Nbar(\lob,\zob)
= \pi\,\thob^2
\int\!dz\;\frac{dV}{dz\,d\Omega}(z)\; w_z(z,\zob)\;
\nrnd(z),
\;}
\label{eq:P1_1d}
\end{equation}
with $\nrnd(z)$ the mean eligible richness density of
eq.~\eqref{eq:blam} --- one smooth single-peaked 1-D integral;
the $\theta$ and $\lambda$ axes disappear analytically. Read it
as a count: (aperture area) $\times$ (eligible galaxies per unit
volume) $\times$ (photo-$z$ window), summed down the cylinder.
(The production code currently applies the $\theta_{\rm excl}(z)$
lower cut to $\Nbar$ as well; whether percolation physically
excludes the Poisson-channel counts inside the exclusion ring is
a modelling choice worth an explicit decision --- the difference
is a ring-sliver correction of relative size
$O(\delta z_{\rm excl}/\sigma_z)\sim 10^{-2}$ of the ring
contribution, itself a small fraction of the total. The
validation of \S\ref{sec:validation} applies the production
convention, so the comparison is convention-for-convention.)

\subsection{The clustered kernel \texorpdfstring{$\Klss$}{Klss}, near zone: integrate on the sphere}
\label{sec:near}

Within $|\pi|\le\pi_s$ the weights are frozen
(\S\ref{sec:frozen}) and the integrand of the clustered kernel is
(counting kernel at transverse offset) $\times$ (spherical
clustering weight with a spherical hole). Everything is a
function of the 3-D separation $r$ and the direction cosine
$\mu = \pi/r$; with $r^2=\pi^2+\rp^2$ and
$d\pi\,\rp\,d\rp = r^2\,dr\,d\mu$, the near zone
$\{r>\Rex,\ \rp\le2\Rex,\ |\pi|\le\pi_s\}$ becomes
\begin{equation}
\boxed{\;
K_X^{\rm near}
= 2\pi \int_{\Rex}^{r_{\max}} dr\; r^2\,\ximm(r)\; G_X(r),
\qquad
G_X(r) = 2\!\int_{\mu_{\rm lo}(r)}^{\mu_{\rm hi}(r)}\! d\mu\;
\Ncal_X\Big(\tfrac{r}{\Rex}\sqrt{1-\mu^2}\Big)\,\bar w_z(r\mu),
\;}
\label{eq:master_near}
\end{equation}
with $X\in\{b,\sigma\}$ ($\Ncal_\sigma \equiv \sigma\Ncal_{\rm lss}$ for
$\KlssInf$), $\mu_{\rm lo}(r) = \sqrt{\max(0,\,1-(2\Rex/r)^2)}$
(the cylinder wall), $\mu_{\rm hi}(r) = \min(1,\,\pi_s/r)$ (the
near/far split), $r_{\max}=\sqrt{\pi_s^2+4\Rex^2}$, and
$\bar w_z$ the two-sided average of the photo-$z$ kernel (unity
to $O((\pi_s/\sigma_\chi)^2)$; kept exactly at negligible cost).

Everything sharp has become smooth:
\begin{itemize}
\item the exclusion is the \emph{exact lower limit} $r=\Rex$, not
      a discontinuity inside the domain;
\item $\ximm(r)$ is a locally power-law decay, ideal for a
      log-$r$ grid;
\item $G_X(r)$ is smooth except for a derivative kink at
      $r=2\Rex$ where the cylinder wall engages --- split the
      $r$-interval there and both pieces are analytic.
\end{itemize}
Figure~\ref{fig:radial} shows the radial integrand: a featureless
power-law decay. The twin exclusion peaks and ring plateau of the
original $z$-integrand are absent --- they were never physical
structure, only the shadow of the excluded ball sliced along $z$.

\begin{figure}[!htbp]
\centering
\includegraphics[width=0.72\textwidth]{figs/pedag_frozen_radial.png}
\caption{The near-zone radial integrand
$r^2\,\ximm(r)\,G_X(r)$ of eq.~\eqref{eq:master_near} at
$(\lob,\zob)=(20,0.5)$. The exclusion radius is the hard lower
limit (dotted); the only non-analytic point is the derivative
kink at $r=2\Rex$ (dashed) where the cylinder wall $\rp\le2\Rex$
engages. Compare with the twin-peak + plateau $z$-integrand of
the original coordinates (\texttt{richness\_selection.tex},
Fig.~\texttt{fz}): identical information, hostile presentation.}
\label{fig:radial}
\end{figure}

\subsection{The clustered kernel, far zone: Limber factorisation with moment corrections}
\label{sec:far}

For $|\pi|>\pi_s \gg 2\Rex$ the separation is dominated by the
line of sight, $r = |\pi|\sqrt{1+\rp^2/\pi^2}$, and $\ximm$
Taylor-expands under the $\rp$-integral:
\begin{equation}
\ximm(r) \simeq \ximm(|\pi|) +
\frac{\rp^2}{2|\pi|}\,\ximm'(|\pi|) + \dots
\end{equation}
The $\rp$-integrals are then precisely the closed-form moments of
\S\ref{sec:moments}. Defining
\begin{equation}
A_X = \int_0^{2\Rex}\!\rp\, \Ncal_X\,d\rp
    = \Rex^2\!\int_0^2\! x\, \Ncal_X\,dx,
\qquad
B_X = \int_0^{2\Rex}\!\rp^3\, \Ncal_X\,d\rp
    = \Rex^4\!\int_0^2\! x^3 \Ncal_X\,dx,
\end{equation}
(for $\Klss$ both are closed by the total-donation theorem:
$A_b = \Rex^2\,\nlss/2$ and
$B_b = \Rex^4\,\avg{b\lambda(1+x_\lambda^2)}/4$; for $\KlssInf$
the $\sigma$-weighted moments $A_\sigma$, $B_\sigma$ are one-time
1-D integrals of the universal kernel), the far zone is the 1-D
Limber-like integral
\begin{equation}
\boxed{\;
K_X^{\rm far}
= 2\pi\sum_{\pm}\int_{\pi_s}^{\pi_{\max}^{\pm}} d\pi\;
\Big(1\pm\tfrac{\pi}{\chi_o}\Big)^{\!2}\, w_z\big(z(\pm\pi)\big)\,
s_{\rm w}\big(z(\pm\pi)\big)
\left[ A_X\,\ximm(\pi) + \frac{B_X}{2\pi}\,\ximm'(\pi)\right],
\;}
\label{eq:master_far}
\end{equation}
where $s_{\rm w}(z) = \nlss(z)/\nlss(\zob)$ is the (slow)
amplitude drift of the frozen weights, evaluated on a coarse
$z$-spline --- the far zone keeps the full $z$-dependence because
it is 1-D and cheap. The $(\chi/\chi_o)^2$ factor carries
\emph{two} powers of the distance ratio: one from the measure of
eq.~\eqref{eq:measure}, and one from the aperture-moment dilation
--- the kernel argument is $x = \theta/\thob =
(\rp/\Rex)\sqrt{\chi_o/\chi}$, so the $\rp$-moments $A_X$, $B_X$
computed at frozen scale pick up $\chi/\chi_o$ per two powers of
$\rp$. In the near zone both powers are $1+O(\pi/\chi_o)$ with
odd-parity cancellation; in the far wings ($|\pi|$ up to
$\sigma_\chi$, $\chi/\chi_o$ up to $\pm10\%$) the squared factor
is required --- with a single power the $\Nbar$ wings, which
carry full weight with no $\ximm$ suppression, err at the
$0.5\%$ level; with the square they return to $\lesssim0.015\%$.
The expansion parameter is $(2\Rex/\pi_s)^2 = 10^{-2}$ at
$\pi_s = 20\Rex$; with the first correction retained the residual
is $O((2\Rex/\pi_s)^4)\sim 10^{-4}$.

Figure~\ref{fig:farzone} shows the far-zone integrand: the BAO
bump at $\approx 105\,\hMpc$ sits well inside the photo-$z$
window ($\sigma_\chi\approx 270\,\hMpc$), which is why the far
zone cannot be truncated early --- distant correlated structure
still donates eligible galaxies as long as $w_z$ lets them pass
the photo-$z$ cut.

\begin{figure}[!htbp]
\centering
\includegraphics[width=0.72\textwidth]{figs/pedag_frozen_farzone.png}
\caption{Far-zone line-of-sight integrand (background side,
log-$\pi$ measure) at $(\lob,\zob)=(20,0.5)$, normalised to its
value at $\pi_s=20\Rex$. The BAO feature at $\approx 105\,\hMpc$
(dotted) lies inside the parabolic photo-$z$ window (dashed): the
far tail carries real weight and the full $w_z(z)\,\nlss(z)$
dependence is retained --- trivially, since the integral is 1-D.}
\label{fig:farzone}
\end{figure}

\subsection{The LSS-asymptotic kernel \texorpdfstring{$\KlssInf$}{Khm-inf}, for free}
\label{sec:I1}

$\KlssInf$ is the clustered kernel weighted toward large $\theta$
--- the part of the cylinder-galaxy excess that lands where the
sigmoid has saturated ($\sigma\to1$) and the bias has settled
onto its large-scale plateau. It needs no new derivation:
substitute $\Ncal_\sigma=\sigma\Ncal_{\rm lss}$ for $\Ncal_{\rm lss}$ in
eqs.~\eqref{eq:master_near} and \eqref{eq:master_far} ($A_b,B_b
\to A_\sigma,B_\sigma$) and every other factor --- $\ximm$,
$\bar w_z$, $w_z$, $s_{\rm w}(z)$, the geometry --- is shared,
unchanged. Assembled:
\begin{equation}
\boxed{\;
\Klss = K_b^{\rm near}+K_b^{\rm far},
\qquad
\KlssInf = K_\sigma^{\rm near}+K_\sigma^{\rm far}.
\;}
\label{eq:K_assembled}
\end{equation}
The entire cost of $\KlssInf$ beyond $\Klss$ is the cheap
$\sigma$-weighted moments and one more $(r,\mu)$/$\pi$ pass
reusing the same $\ximm$, $\bar w_z$, $w_z$ evaluations. Because
the two kernels share one pipeline, whatever approximation error
enters $\Klss$ enters $\KlssInf$ nearly identically --- the
correlated errors that keep the $\Klss-\KlssInf$ denominator of
\S\ref{sec:closure} well-conditioned
(\S\ref{sec:bsel_validation}).

\section{From counts to the selection bias}
\label{sec:closure}

The counts close the selection bias by budget balance: the
observed inflation must equal Poisson count plus LSS count, where
the LSS galaxies are counted with the \emph{target's} own bias
--- completing the two-tracer correlation of
eq.~\eqref{eq:two_tracer}. For a richness-selected target that
bias is the unknown $\bsel(\theta)$, split by the sigmoid ansatz
into two plateaus,
$\bsel(\theta) = \bsmall[1-\sigma(\theta)]+\blarge\sigma(\theta)$:
\begin{equation}
\lob-\ltr
\;=\; \underbrace{\Nbar}_{\text{Poisson}}
\;+\; \underbrace{\blarge\,\KlssInf + \bsmall\,(\Klss-\KlssInf)}_{\text{LSS,
counted with the target's bias}}.
\label{eq:budget}
\end{equation}
Nothing in this section is new to the reformulation --- it is the
production algebra of \texttt{richness\_selection.tex},
\S``Bias pipeline'', reproduced so this note is self-contained;
the only change is that the counts on the right are now the
reformulated quantities of \S\ref{sec:counts}.

The remaining inputs are exact and shared with production. The
target's own bias for a \emph{mass-selected} (randomly-selected)
sample at the same $(\lob,\zob)$ --- a 1-D $\ln M$ integral at
fixed $\zob$, nothing to freeze ---
\begin{equation}
\btar(\lob,\zob)
= \frac{\int d\ln M\;M\,n(M,\zob)\,b(M,\zob)\,P(\lob\mid M,\zob)}
       {\int d\ln M\;M\,n(M,\zob)\,P(\lob\mid M,\zob)}
\label{eq:btar}
\end{equation}
anchors the random-aperture expectation: substituting
$\bsel\to\btar$ (a constant, which pulls out of the mass
integral) in eq.~\eqref{eq:budget} gives the contamination a
random mass-matched target would suffer,
\begin{equation}
\Nrnd(\lob,\zob) \;=\; \Nbar \;+\; \Nlss,
\qquad
\Nlss \;\equiv\; \btar\,\Klss
\label{eq:Nrnd}
\end{equation}
--- Poisson plus LSS-clustered contamination, the same
Poisson$+$clustering split as
$\Sigma=\Sigma^{\rm 1h}+\Sigma^{\rm prj}$ one level up,
denominated in galaxy counts (production:
eq.~\texttt{dprj\_rnd}). The large-scale plateau follows from
the Buzzard-calibrated empirical relation (production:
eq.~\texttt{b\_infty}):
\begin{equation}
\blarge(\lob,\ltr,\zob) = \btar(\lob,\zob)
  \big[1+0.13\,\delta^{\rm prj}\big],
\qquad
\delta^{\rm prj}=\frac{\lob-\ltr}{\Nrnd}-1,
\label{eq:blarge_frozen}
\end{equation}
with $\delta^{\rm prj}$ the fractional excess contamination of
the actual $(\lob,\ltr)$ pair over the random expectation. The
small-$\theta$ plateau then closes the budget,
eq.~\eqref{eq:budget}, as the remaining linear unknown
(production: eq.~\texttt{b\_small}):
\begin{equation}
\boxed{\;
\bsmall(\lob,\ltr,\zob) =
\frac{(\lob-\ltr)-\Nbar-\blarge\,\KlssInf}{\Klss - \KlssInf}.
\;}
\label{eq:bsmall_frozen}
\end{equation}
This numerator near-cancels at small $\dprj=\lob-\ltr$
(\S\ref{sec:bsel_validation}) --- the demanding test of any
approximation to the counts. With $\bsmall$, $\blarge$ in hand,
the sigmoid ansatz returns $\bsel(\lob,\ltr,\zob,\theta)$;
$\ltr$-marginalisation
(\texttt{richness\_selection.tex}, eq.~\texttt{b\_marg\_lt}) is
unaffected by anything in this note.

\section{Validation}
\label{sec:validation}

\subsection{The three counts against \texorpdfstring{\texttt{scipy.quad}}{scipy.quad}}
\label{sec:quad_validation}

The complete reformulated computation of $(\KlssInf, \Klss)$ at one
$(\lob,\zob)$:
\begin{enumerate}
\item contract the eligible $(M,\lambda)$ population into
      $\Ncal_{\rm lss}(x)$ (eq.~\ref{eq:N_def}); set
      $\Ncal_\sigma = \sigma\Ncal_{\rm lss}$;
\item near zone by eq.~\eqref{eq:master_near}: log-GL in $r$ on
      $[\Rex,2\Rex]\cup[2\Rex,r_{\max}]$, GL in $\mu$;
\item far zone by eq.~\eqref{eq:master_far}: log-GL in $\pi$ per
      side, moments $A_X$, $B_X$ precomputed, $\nlss(z)$ drift on
      a coarse spline;
\item $K_X = K_X^{\rm near} + K_X^{\rm far}$; $\Nbar$ by
      eq.~\eqref{eq:P1_1d}.
\end{enumerate}
Validation against \texttt{scipy.integrate.quad} on the original
$(z,\theta,M,\lambda)$ formulation (adaptive in $z$ at
$\epsilon_{\rm rel}=10^{-6}$, matched inner grids), for both the
production grid code (\texttt{SelBias.bias\_precompute}, current
repository defaults $N_{z,\rm bias}=48$, $N_\theta=10$) and the
frozen reformulation:
\begin{center}
\small
\begin{tabular}{r r l r r r}
\toprule
$\lob$ & $\zob$ & source & $\Nbar$ err & $\KlssInf$ err & $\Klss$ err \\
\midrule
$20.0$ & $0.500$ & production & $0.018\%$ & $0.004\%$ & $0.002\%$ \\
       &         & frozen     & $0.006\%$ & $0.037\%$ & $0.037\%$ \\
$52.5$ & $0.425$ & production & $0.007\%$ & $0.020\%$ & $0.021\%$ \\
       &         & frozen     & $0.015\%$ & $0.030\%$ & $0.019\%$ \\
\bottomrule
\end{tabular}
\end{center}
Both are comfortably below the $0.1\%$ pipeline tolerance; the
frozen path replaces the entire $(z,M,\lambda,\theta)$ 4-D
quadrature by one $(M,\lambda)$ contraction, one smooth 2-D
$(r,\mu)$ block, and two 1-D integrals. (For $\Nbar$ the frozen
path applies the same exclusion convention as production --- the
ball sliver enters as the incomplete moment
$M_0(2)-M_0(x_{\rm excl}(\pi))$ with
$M_0(y)=\int_0^y x\,\Ncal_{\rm rnd}(x)\,dx$ --- so the comparison is
convention-for-convention.)

\subsection{End-to-end: the selection-bias profile and the lensing observable}
\label{sec:bsel_validation}

The counts are inputs; the pipeline output is
$\bsel(\theta)$ assembled through \S\ref{sec:closure}. This is
the demanding test: the small-$\theta$ plateau
eq.~\eqref{eq:bsmall_frozen} is a near-cancellation at small
$\dprj=\lob-\ltr$, so fractional errors on the counts are
\emph{amplified} in the numerator. To isolate the reformulation,
both paths share the same downstream algebra
(\texttt{SelBias.bias\_from\_precomp}) and the same $\btar$;
only $(\Nbar, \KlssInf, \Klss)$ differ --- production grid values
vs.\ frozen assembly.

Both figures below swap only \texttt{SelBias.bias\_precompute}
between production and frozen; every other line of the pipeline
(NFW lookup, $\theta$-grid, aperture-weighting, plateau algebra)
is byte-identical between the two curves in each panel. The
second figure propagates the comparison one step further, into
the lensing observable $\avg{\Delta\Sigma^{\rm prj}(R)}$
(\texttt{delta\_sigma\_prj.py}), which uses $\bsel(\theta)$ as an
input kernel inside the projected-NFW integral.

\begin{figure}[!htbp]
\centering
\includegraphics[width=0.92\textwidth]{figs/pedag_frozen_vs_prod_bsel.png}
\caption{$\bsel(\theta\mid\lob,\ltr,\zob)$ at
$(\lob,\zob)=(20,0.5)$ for $\dprj\in\{0,3,6,10,15\}$: production
(dashed) vs.\ frozen reformulation (solid, visually coincident).
\emph{Bottom:} residual normalised by
$\max_\theta|b^{\rm prod}|$ per $\dprj$; largest at $\dprj=0$
(the sign-changing profile, where the small-$\theta$ plateau is a
near-cancellation and amplifies input errors) but still
$\lesssim 0.05\%$ of the profile scale.}
\label{fig:bsel_vs_prod}
\end{figure}

\begin{figure}[!htbp]
\centering
\includegraphics[width=0.8\textwidth]{figs/pedag_frozen_vs_prod_delta_sigma.png}
\caption{$\avg{\Delta\Sigma^{\rm prj}(R\mid\lob,\zob)}$ (cl+LSS
piece) at the same $(\lob,\zob,\dprj)$ grid, production
(dashed, circles) vs.\ frozen (solid, squares); curves overlap at
plot resolution. \emph{Bottom:} fractional residual
$\Delta\Sigma^{\rm froz}/\Delta\Sigma^{\rm prod}-1$,
$\lesssim0.1\%$ throughout, confirming the $\bsel(\theta)$
agreement survives being folded through the $\theta$-integral
against the miscentered-NFW kernel. Points within $5\%$ of the
$\dprj=0$ curve's sign crossing (top panel;
$\Delta\Sigma^{\rm prj}(R)$ itself passes through zero there) are
dropped from the residual: the fractional ratio diverges at the
crossing even though the absolute difference stays at the same
sub-percent level as everywhere else.}
\label{fig:dsigma_vs_prod}
\end{figure}

Figure~\ref{fig:bsel_vs_prod} shows the result: the frozen
counts reproduce the production $\bsel(\theta)$ to well below
$0.1\%$ of the profile scale across all $\theta$ and all
contamination levels, including the sign-changing $\dprj=0$
profile ($0.07\%$ worst case at $\dprj=2$, measured in the
\S\ref{sec:quad_validation} run). The plateau-cancellation
amplification is visible but mild --- count errors of
$\sim0.04\%$ surface as at most $\sim0.07\%$ on the profile ---
because the $\Klss-\KlssInf$ denominator and the $\Nbar$
subtraction are hit by \emph{correlated} errors (near and far
zones are shared between the two kernels, \S\ref{sec:I1}), which
largely cancel in the ratio.

\subsection{Error budget}
\begin{center}
\small
\begin{tabular}{l l l}
\toprule
approximation & order & size \\
\midrule
$\rp$ variable (eq.~\ref{eq:rp_def}) & exact & $0$ \\
$\theta_{\max}$ cut vs.\ $f_A$ support & exact for $x_\lambda\le1$ & $\sim0$ \\
total-donation theorem (eq.~\ref{eq:moments}) & exact & $0$ \\
$\KlssInf$ kernel $=\sigma\Ncal_{\rm lss}$ & exact & $0$ \\
frozen weights, near zone & odd part cancels; $(\Delta z\,\partial_z)^2$ & $\sim10^{-4}$ \\
$(1+\pi/\chi_o)^2$ frozen, near zone & odd part cancels; $(\pi_s/\chi_o)^2$ & $\sim10^{-4}$ \\
$(1+\pi/\chi_o)^2$, far zone & kept exactly (both powers) & $0$ \\
far-zone moment expansion & $(2\Rex/\pi_s)^4$ & $\sim10^{-4}$ \\
\emph{assembled total (measured)} & & $0.6$--$4\times10^{-4}$ \\
\emph{end-to-end $\bsel(\theta)$ (measured)} & plateau cancellation, correlated-error relief & $\le 7\times10^{-4}$ \\
\bottomrule
\end{tabular}
\end{center}

\subsection{Stress test: MOR redshift-evolution slope}
\label{sec:mor_stress}

The frozen-weights argument of \S\ref{sec:frozen} rests on the MOR
$P(\lambda\mid M,z)$ evolving slowly across the exclusion-affected
zone. The MOR's own redshift-evolution rate is set by a single
parameter, the satellite-count slope $\epsilon$ in
$l_{\rm sat}(M,z)\propto\big[(1+z)/(1+z_{\rm pivot})\big]^{\epsilon}$
(\texttt{MOR.epsilon}, fiducial default $\epsilon=0$); it is the one
knob that speaks directly to the frozen approximation's core
assumption, so it is the right dial to stress rather than assume.

Because the question that matters operationally is not ``how close
is frozen to \texttt{scipy.quad}'' but ``how much would swapping
frozen in for the deployed pipeline change the answer,'' the
comparison here is \emph{frozen against production directly} ---
both evaluated on the identical $(\text{cosmology},\text{MOR})$
stack at each $\epsilon$ --- rather than each against \texttt{quad}
separately (the two vs.-quad errors are uncorrelated approximations
of different origin and need not partially cancel, so a direct
frozen-vs-production comparison is the stronger test):
\begin{center}
\small
\begin{tabular}{r r r r r}
\toprule
$\lob$ & $\epsilon$ & $|\Delta P_1/P_1|$ & $|\Delta\Klss/\Klss|$ & $|\Delta\KlssInf/\KlssInf|$ \\
\midrule
$20.0$   & $-2.0$ & $0.0004\%$ & $0.0312\%$ & $0.0294\%$ \\
$20.0$   & $0.0$  & $0.0055\%$ & $0.0297\%$ & $0.0274\%$ \\
$20.0$   & $+2.0$ & $0.0103\%$ & $0.0280\%$ & $0.0246\%$ \\
$52.5$   & $-2.0$ & $0.0064\%$ & $0.0175\%$ & $0.0036\%$ \\
$52.5$   & $0.0$  & $0.0002\%$ & $0.0183\%$ & $0.0011\%$ \\
$52.5$   & $+2.0$ & $0.0028\%$ & $0.0185\%$ & $0.0059\%$ \\
$130.0$  & $-2.0$ & $0.0734\%$ & $0.0025\%$ & $0.0134\%$ \\
$130.0$  & $0.0$  & $0.0483\%$ & $0.0052\%$ & $0.0050\%$ \\
$130.0$  & $+2.0$ & $0.0275\%$ & $0.0137\%$ & $0.0225\%$ \\
\bottomrule
\end{tabular}
\end{center}
at $\zob=0.60$ ($\Delta X \equiv X^{\rm frozen}-X^{\rm prod}$, both
on the same MOR); full 9-point $\epsilon\in[-2,2]$ sweep in
\texttt{validations/cache/frozen\_mor\_epsilon\_stress.csv}, produced
by \texttt{validations/frozen\_mor\_epsilon\_stress.py}. Two
findings:
\begin{itemize}
\item \textbf{The number.} Worst case over the full sweep is
      $0.073\%$ ($P_1$ at $\lob=130$, $\epsilon=-2$) --- inside the
      $0.1\%$ pipeline tolerance, but with less margin than the
      vs.-\texttt{quad} comparisons of \S\ref{sec:quad_validation}
      suggested in isolation.
\item \textbf{The slope.} Only $P_1$ at $\lob=130$ shows a clear
      \emph{monotonic} trend with $\epsilon$ (0.028\% at
      $\epsilon=+2$ rising to $0.073\%$ at $\epsilon=-2$); every
      other $(\lob,X)$ combination is flat or non-monotonic across
      the sweep, at 3--10$\times$ smaller amplitude. The trend tracks
      $P_1$ specifically because $\Nbar$'s frozen amplitude-drift
      spline $s_{\rm w}(z)$ (\S\ref{sec:far}) is the one piece of
      the reformulation that carries the MOR's $z$-dependence
      explicitly beyond the near zone; $\Klss,\KlssInf$ do not show
      it because production's own $N_{z,\rm bias}=48$ grid error is
      already the dominant term there (\S\ref{sec:relation}).
\end{itemize}
This is a stress test, not a physical regime: the DES Y1 NC+3x2pt
calibration \texttt{MOR} defaults to (and is centred near)
$\epsilon=0$ (\texttt{mor.py}), far inside the flat part of this
sweep. The result that matters is qualitative: nothing here breaks
the frozen reformulation even at $\epsilon=\pm2$, several times the
plausible astrophysical range.

\section{Relation to the production recipe}
\label{sec:relation}

The production ring+outer GL split
(\texttt{richness\_selection.tex}, \S``The $z$-axis'') and this
reformulation attack the same integrand from opposite directions:
the former keeps the original coordinates and engineers the grid
around the sharp features; the latter changes coordinates so the
sharp features never form. They agree at the
$\lesssim 4\times10^{-4}$ level at the fiducial
$(\lob,\zob)=(20,0.5)$ point, and stay $\le 7.3\times10^{-4}$
under the $\epsilon\in[-2,2]$ MOR-slope stress test of
\S\ref{sec:mor_stress} at $\zob=0.60$ --- the two recipes' residual
disagreement is not sensitive to how fast the MOR itself evolves in
$z$. The reformulation additionally
buys: the exact $\sigma\Ncal_{\rm lss}$ relation for the $\KlssInf$ kernel
(halving the kernel work), the total-donation theorem
(collapsing $\Nbar$ and the far-zone $\Klss$ amplitude), and a
2-D-instead-of-4-D numerical core.

Reference implementation:
\texttt{validations/frozen\_kernels.py} (kernels, assembly, and
the \S\ref{sec:kernels}--\ref{sec:counts} figures;
code symbols follow the production names of the dictionary in
\S\ref{sec:question});
\texttt{docs/make\_validation\_plots\_frozen\_physics.py} (the
end-to-end comparison figures of \S\ref{sec:bsel_validation});
frozen-ring validation table of \S\ref{sec:frozen} from the ring
freeze test. Production \texttt{SelBias.\_P\_operator} is
unchanged by this note.
\fi

\end{document}
